\documentclass[preprint,12pt,authoryear]{elsarticle}

\usepackage{amsmath,amssymb,amsthm}
\usepackage{booktabs}
\usepackage{multirow}
\usepackage{xcolor}
\usepackage{graphicx}
\usepackage{tikz}
\usetikzlibrary{arrows.meta,positioning,fit,backgrounds,calc,shapes.geometric}
\usepackage{hyperref}

\newtheorem{proposition}{Proposition}
\newtheorem{remark}{Remark}
\newcommand{\placeholder}[1]{\textcolor{red}{[#1]}}
\newcommand{\norm}[1]{\left\lVert #1 \right\rVert}

\journal{Transportation Research Part E: Logistics and Transportation Review}

\usepackage{xeCJK}
\setCJKmainfont[AutoFakeBold=2.5,AutoFakeSlant=0.2]{Noto Serif CJK SC}
\setCJKsansfont{Noto Sans CJK SC}
\setCJKmonofont{Noto Sans Mono CJK SC}
\definecolor{zhc}{RGB}{25,65,145}
\newcommand{\zh}[1]{\par\vspace{2pt}{\color{zhc}\small #1\par}\vspace{6pt}}
\newcommand{\zhitem}[1]{\par{\color{zhc}#1}}
\DeclareRobustCommand{\zhcap}[1]{{\color{zhc}#1}}
\renewcommand{\topfraction}{0.9}
\renewcommand{\bottomfraction}{0.8}
\renewcommand{\textfraction}{0.07}
\renewcommand{\floatpagefraction}{0.6}
\setcounter{topnumber}{3}
\setcounter{totalnumber}{5}

\begin{document}

\begin{frontmatter}

\title{Deep Reinforcement Learning with Few Exact Solves for the Carrier--Vehicle Traveling Salesman Problem\\[6pt]{\large 基于少量精确求解的深度强化学习求解载体—无人机旅行商问题}}

\author{\placeholder{Author names, affiliations and corresponding author}}

\begin{abstract}
In the carrier--vehicle traveling salesman problem (CVTSP), a carrier launches and recovers a drone at continuous points while the drone visits a set of targets. Evaluating one visiting order requires solving a second-order cone program (SOCP), so training a deep reinforcement learning (DRL) policy on exact rewards for every sampled route consumes hundreds of millions of solver calls. We show that solving one order already gives a feasible solution for every other order: keep the take-off and landing points of each target and let the carrier drive between them in the new order. The cost of this solution, read from a table without a solver, is therefore an upper bound on the optimal cost of any order, and it is exact for the order that was solved. We build a training recipe around this bound: it gives a reward to every sampled route, guides solver-free local search and ranks candidate expert routes, while exact solves are reserved for a few routes per instance and for every imitation label. With the same 3.84 million exact solves, POMO trained on exact rewards is on average 3.8\% worse per instance than our policy on 529 literature benchmark instances. A POMO baseline trained with 299 million solves, 78 times as many as ours in training, remains better than our policy, by 0.21\% on the same instances and 1.13\% on synthetic test instances. At equal computing time, iterated local search and variable neighbourhood search are worse than our policy on these 529 instances, but with ten times the time they find better solutions than both learned policies. The recipe is therefore a way to obtain a learned CVTSP solver with a small solver budget, not a better solver.

\zh{在载体—无人机旅行商问题（CVTSP）中，载体在连续平面上的点释放和回收无人机，无人机依次访问一组目标。评估一个访问顺序需要求解一个二阶锥规划（SOCP），因此若对每条采样路线都用精确奖励训练深度强化学习（DRL）策略，需要数亿次求解器调用。我们证明，求解一个顺序就已经为其他每个顺序给出了一个可行解：保留每个目标的起飞点和降落点，让载体按新的顺序在这些点之间行驶。这个解的代价可以从一张表中读出而无需求解器，因此它是任意顺序最优代价的上界，并且对已求解的那个顺序是精确的。我们围绕这一上界构建训练方案：它为每条采样路线提供奖励、引导无需求解器的局部搜索并对候选专家路线排序，而精确求解只用于每个实例的少数路线以及每个模仿标签。在同样使用 384 万次精确求解时，以精确奖励训练的 POMO 在 529 个文献基准实例上逐实例平均比我们的策略差 3.8\%。一个用 2.99 亿次求解训练的 POMO 基线（训练求解次数是我们的 78 倍）仍优于我们的策略，在同一批实例上差距为 0.21\%，在合成测试实例上为 1.13\%。在同等计算时间下，迭代局部搜索和变邻域搜索在这 529 个实例上不如我们的策略，但给予十倍时间时，它们能得到比两个学习型策略都更好的解。因此，本方案是以较小的求解器预算获得一个学习型 CVTSP 求解器的方法，而不是一个更好的求解器。}
\end{abstract}

\begin{keyword}
Carrier--vehicle routing \sep Drone logistics \sep Deep reinforcement learning \sep Second-order cone programming \sep Certified upper bound
\end{keyword}

\end{frontmatter}

\section{Introduction\quad 引言}
\label{sec:intro}

A drone with a short flight range can serve targets far beyond that range if a long-range platform carries it. A ship can launch an unmanned aerial vehicle (UAV) towards an offshore installation, keep moving, and recover the UAV further along its course. The carrier--vehicle traveling salesman problem (CVTSP) \citep{garone2011} models this operation. A carrier leaves a depot, releases a drone before each target, recovers it after the visit and returns to the depot once all targets have been served. Each release, visit and recovery is a sortie, and the carrier's travel from the landing point of one sortie to the take-off point of the next is a carrier leg. The planner chooses the visiting order and the take-off and landing point of each sortie in order to minimise mission completion time.

\zh{航程较短的无人机若由远程平台携带，就能服务远超其航程的目标。船舶可以向海上设施发射一架无人机（UAV），继续航行，并在航线更远处回收它。载体—无人机旅行商问题（CVTSP）\citep{garone2011} 刻画的就是这种作业。载体从基地出发，在每个目标之前释放无人机，在访问之后回收它，并在服务完所有目标后返回基地。每一次释放、访问和回收称为一个架次；载体从一个架次的降落点行驶到下一个架次起飞点的这段路程称为一个载体航段。规划者需要选择访问顺序以及每个架次的起飞点和降落点，使任务完成时间最短。}

The problem combines two kinds of decisions: the visiting order, which is a permutation, and the take-off and landing points of the drone, which lie in the continuous plane. Once the order is fixed, the remaining problem is convex and can be written as an SOCP \citep{lobo1998}; when the order is also a decision, the problem becomes a mixed-integer SOCP. Exact methods solve it only for small instances \citep{gambella2018,erdogan2021}; the largest instances proved optimal so far have 25 targets, plus one with 35 \citep{vieira2026}. Larger instances are solved by heuristics, such as the two-stage iterative heuristic of \citet{li2025}, which compute every instance from scratch.

\zh{该问题结合了两类决策：访问顺序是一个排列，而无人机的起飞点和降落点位于连续平面上。顺序固定后，剩余问题是凸的，可写成 SOCP \citep{lobo1998}；当顺序也是决策变量时，问题成为混合整数 SOCP。精确方法只能求解小规模实例 \citep{gambella2018,erdogan2021}；目前被证明最优的最大实例有 25 个目标，另有一个 35 个目标的实例 \citep{vieira2026}。更大的实例依靠启发式求解，例如 \citet{li2025} 的两阶段迭代启发式，它们对每个实例都从头计算。}

Deep reinforcement learning has become a practical alternative for routing problems when solutions must be produced quickly. A trained policy produces a route in a single forward pass, and methods such as POMO \citep{kwon2020} reach small gaps on the traveling salesman problem. These methods are trained by sampling millions of routes and scoring each one, which is cheap when the cost of a route is a sum of known edge lengths. In the CVTSP it is not: the carrier legs depend on where the drone takes off and lands, and these points are optimised jointly for the whole route, so scoring one route means solving an SOCP. Our reference baseline, POMO trained on an exact solver reward for every sampled route, made about $3 \times 10^{8}$ solver calls over 200 epochs. Within the range of speeds and endurance seen in training, such a policy can be reused without retraining, because these parameters enter the network as inputs. A new variant of the problem, such as one with additional constraints or a different spatial distribution of targets, may require training from scratch. The solver budget of that training then decides whether building a learned solver is worthwhile at all.

\zh{当需要快速给出解时，深度强化学习已成为路径问题的一种实用替代方法。训练好的策略一次前向计算即可生成路线，POMO \citep{kwon2020} 等方法在旅行商问题上已达到很小的差距。这些方法在训练中要采样数百万条路线并逐条打分；当路线代价是已知边长之和时，这一步代价很低。CVTSP 不是这样：载体航段取决于无人机在哪里起飞和降落，而这些点是针对整条路线联合优化的，因此给一条路线打分就要求解一个 SOCP。我们的参照基线，即对每条采样路线都使用精确求解器奖励训练的 POMO，在 200 个训练轮次中进行了约 $3 \times 10^{8}$ 次求解器调用。在训练中见过的速度和续航范围内，这样的策略无需重新训练即可复用，因为这些参数作为输入进入网络。而问题的新变体，例如增加了约束或目标的空间分布不同，可能需要从头训练。此时，训练所需的求解器预算决定了构建学习型求解器是否值得。}

This paper asks whether most of those calls are necessary. The answer rests on one structural observation. In the fixed-order SOCP, the constraints on a sortie concern one target only, and the visiting order enters only through the carrier legs between sorties. If the sortie points from one optimal solution are kept and the order is changed, every sortie remains feasible and only the carrier legs change. The resulting value is an upper bound on the optimal cost of the new order, and it can be computed by adding up entries of a table. One exact solve thus says something about every route of the instance, including the routes that will never be solved.

\zh{本文要回答的问题是：这些调用中的大多数是否必要。答案建立在一个结构性观察之上。在定序 SOCP 中，每个架次的约束只涉及一个目标，访问顺序只通过架次之间的载体航段进入问题。如果保留一个最优解中的架次点而改变顺序，每个架次仍然可行，只有载体航段发生变化。由此得到的值是新顺序最优代价的一个上界，并且可以通过对一张表中的元素求和得到。因此，一次精确求解就能提供该实例中每条路线的信息，包括那些永远不会被求解的路线。}

\citet{erdogan2021} already use the idea of keeping sortie points fixed. In their iterated local search, a neighbouring order keeps the take-off and landing points of the current solution, so its cost can be computed without a new solve. What we have not found is the use of this idea to train a DRL policy. Self-improvement methods, such as the multi-action method of \citet{luttmann2025}, compute the exact cost of every sampled solution, which is cheap in the problems they study. In our recipe, the bound replaces this computation. It gives the reward for routes that are never solved, and it ranks candidate labels, each of which is solved exactly before it is used. Section~\ref{sec:gap} compares our method with the closest CVTSP and learning methods.

\zh{保持架次点不变这一想法，\citet{erdogan2021} 已经用过。在他们的迭代局部搜索中，相邻顺序沿用当前解的起飞点和降落点，因此无需重新求解就能算出其代价。我们没有找到的是用这一想法来训练 DRL 策略。自我改进方法，如 \citet{luttmann2025} 的多动作方法，对每个采样解都直接计算精确代价，这在它们研究的问题中很便宜。在我们的方法中，上界代替了这一计算：它为从未求解的路线提供奖励，并给候选标签排序，每个标签在使用前都经过一次精确求解。第~\ref{sec:gap} 节把我们的方法与最接近的 CVTSP 方法和学习方法做比较。}

The paper makes three contributions. \placeholder{Contributions pending, under discussion with the authors. Every contribution except Proposition~1 is to be tied to an experiment that verifies it. Open: whether the metaheuristic comparison joins contribution~3 or becomes a fourth; what contribution~2 claims before the results of Groups B and D; whether contribution~3 states the pattern by instance size.}
\zh{本文有三项贡献。\placeholder{贡献待定，正在讨论。除命题~1 外，每项贡献都要对应一个验证它的实验。待定问题：元启发式比较并入贡献~3 还是单列为第四项；在 B 组和 D 组结果出来之前，贡献~2 声称什么；贡献~3 是否写出随实例规模变化的规律。}}
\begin{enumerate}
\item We state and prove that one optimal fixed-order solution yields a table whose route sums bound the exact cost of every visiting order from above and match it at the solved order. The proof is short because it uses the structure of the CVTSP directly. The contribution lies in how this bound is put to work. \zhitem{我们陈述并证明：一个最优定序解给出一张表，其路线求和从上方界定每个访问顺序的精确代价，并在已求解的顺序处与之相等。证明很短，因为它直接利用了 CVTSP 的结构。贡献在于如何利用这一上界。}
\item We build a training recipe around the table. Rewards for all sampled routes, local search and the ranking of candidate experts use the table without calling a solver, and exact solves are reserved for the anchor route, for one policy route and for the verification of every imitation label. The recipe uses 2.89 solver calls per training instance on average. \zhitem{我们围绕这张表构建训练方案。所有采样路线的奖励、局部搜索和候选专家的排序都使用该表而不调用求解器，精确求解只用于锚点路线、一条策略路线以及每个模仿标签的验证。该方案平均每个训练实例使用 2.89 次求解器调用。}
\item We compare the recipe with POMO trained on exact rewards for every route, at equal solver budget and at full training, and we report the results by instance size. At equal budget the recipe gives the better policy. The fully trained baseline used 78 times more solver calls in training, or about 29 times more when both runs are counted as run, and it remains better on synthetic instances from the training distribution; we quantify where and by how much. \zhitem{我们将该方案与对每条路线使用精确奖励训练的 POMO 在相同求解器预算和完全训练两种条件下进行比较，并按实例规模报告结果。在相同预算下，本方案得到的策略更好。完全训练的基线在训练中使用的求解器调用是我们的 78 倍，若两边都按实际运行计则约为 29 倍；它在训练分布的合成实例上仍然更好，我们量化了差距出现在哪里、有多大。}
\end{enumerate}

The practical result is a change in what a learned CVTSP solver costs to obtain. Training to epoch 200 on a single workstation took about 28 hours; the baseline run took about 215 hours on a different machine with 128 reward workers. The paper also reports where this comes at a price: the policy does not match the fully trained baseline on instances from its own training distribution, and metaheuristics with longer time budgets obtain better solutions. The contribution is a training method, and we state its practical implications as a comparison of computing budgets rather than as an operational study.

\zh{实际结果是改变了获得一个学习型 CVTSP 求解器的成本。在单台工作站上训练到第 200 轮约需 28 小时；基线的训练在另一台配有 128 个奖励进程的机器上用了约 215 小时。本文也报告了代价所在：该策略在其自身训练分布的实例上不及完全训练的基线，而时间预算更长的元启发式能得到更好的解。本文的贡献是一种训练方法，我们以计算预算比较的方式陈述其实践意义，而不是把它作为一项运营研究。}

Section~\ref{sec:related} reviews related work. Section~\ref{sec:problem} defines the problem and Section~\ref{sec:method} presents the method. Section~\ref{sec:experiments} reports the experiments, Section~\ref{sec:discussion} discusses practical implications and Section~\ref{sec:conclusion} concludes.

\zh{第~\ref{sec:related} 节回顾相关工作。第~\ref{sec:problem} 节定义问题，第~\ref{sec:method} 节介绍方法。第~\ref{sec:experiments} 节报告实验，第~\ref{sec:discussion} 节讨论实践意义，第~\ref{sec:conclusion} 节总结全文。}

\section{Related work\quad 相关工作}
\label{sec:related}

Three strands of work bear on this paper: solution methods for the CVTSP, DRL for routing, and learning under expensive objectives.

\zh{与本文相关的工作有三条脉络：CVTSP 的求解方法、用于路径问题的 DRL，以及昂贵目标下的学习。}

\subsection{The CVTSP and its extensions\quad CVTSP 及其扩展}

\citet{garone2011} studied the problem and proposed a two-step approach that first fixes an order and then optimises the sortie points. \citet{garone2014} studied the variant in which the vehicle visits several targets per sortie. \citet{gambella2018} gave a mixed-integer SOCP formulation and an exact algorithm that solves instances with up to 15 targets; the algorithm ranks visiting orders by a lower bound and solves fixed-order problems in that order, which is an exact-method precedent for screening orders with a cheap bound. \citet{erdogan2021} proposed a formulation with stronger lower bounds and an iterated local search whose neighbourhood moves keep the take-off and landing points of each target, so that a neighbouring order receives a feasible value without a solve. \citet{poikonen2020} studied the mothership and drone routing problem, with instances of up to 200 targets; \citet{erdogan2021} describe it as essentially the same problem as the CVTSP. \citet{fahradyan2020} extended the CVTSP to two carriers that share one vehicle, and \citet{liy2024} studied a joint routing problem with several vessels and UAVs, solved with a tabu search heuristic.

\zh{\citet{garone2011} 研究了该问题，并提出先固定顺序、再优化架次点的两步法。\citet{garone2014} 研究了无人机在一个架次中访问多个目标的变体。\citet{gambella2018} 给出了一个混合整数 SOCP 模型和一个精确算法，可求解最多 15 个目标的实例；该算法按下界给访问顺序排序，并依次求解定序问题，这是精确方法中用廉价下界筛选顺序的先例。\citet{erdogan2021} 提出了下界更强的模型，以及一种迭代局部搜索，其邻域移动保留每个目标的起飞点和降落点，从而无需求解即可为相邻顺序给出一个可行值。\citet{poikonen2020} 研究了母舰—无人机路径问题，其实例最多有 200 个目标；\citet{erdogan2021} 认为它与 CVTSP 本质上是同一问题。\citet{fahradyan2020} 把 CVTSP 扩展到两个载体共用一架无人机的情形，\citet{liy2024} 研究了多艘船舶与多架无人机的联合路径问题，并用禁忌搜索启发式求解。}

Recent work has moved the exact and heuristic frontiers. \citet{li2025} developed exact methods and a two-stage iterative heuristic, and their results supply the reference values used in this paper. \citet{vieira2026} propose a branch-and-cut algorithm that iteratively corrects the error of a linear approximation; it solves all benchmark instances with up to 25 customers to optimality, and one instance with 35 customers. \citet{ren2026} address a vessel with multiple UAVs through a mixed-integer SOCP and a parallel adaptive large neighbourhood search. CVTSP research is one branch of a wider literature on drone-assisted routing \citep{murray2015,agatz2018,otto2018,macrina2020,chung2020}. All of the CVTSP methods above solve each instance from scratch and spend their computing time per instance. A trained policy moves most of that time into training, although at test time it still solves one SOCP per distinct route it proposes.

\zh{近期工作推进了精确方法和启发式方法的前沿。\citet{li2025} 提出了精确方法和一种两阶段迭代启发式，他们的结果提供了本文所用的参考值。\citet{vieira2026} 提出一种分支切割算法，它迭代修正线性近似的误差，求得了最多 25 个客户的全部基准实例的最优解，以及一个 35 个客户实例的最优解。\citet{ren2026} 通过混合整数 SOCP 和并行自适应大邻域搜索处理一艘船携带多架无人机的问题。CVTSP 研究是无人机辅助路径这一更广泛文献的一个分支 \citep{murray2015,agatz2018,otto2018,macrina2020,chung2020}。上述 CVTSP 方法都对每个实例从头求解，把计算时间花在每个实例上。训练好的策略把大部分时间移到训练阶段，不过在测试时，它仍要为提出的每条不同路线求解一个 SOCP。}

Two observations from this literature shape our design. The first is that the fixed-order problem is the unit of exact evaluation: the exact cost of a candidate order requires one SOCP. The stronger heuristics avoid paying it for every candidate, either by keeping sortie points fixed during local search \citep{erdogan2021} or by searching on an approximate matrix \citep{li2025}. The second is that the order and the sortie points are strongly coupled, since a good order depends on where the carrier will be, and the carrier's position depends on the order. Two-stage methods handle this coupling by alternating between the two levels. Our table can be read as one such alternation step, frozen and reused across a whole search.

\zh{这些文献中的两项观察塑造了我们的设计。第一，定序问题是精确评估的基本单位：一个候选顺序的精确代价需要一次 SOCP 求解。较强的启发式方法会避免为每个候选都付出这一代价，要么在局部搜索中保持架次点不变 \citep{erdogan2021}，要么在近似矩阵上搜索 \citep{li2025}。第二，顺序和架次点强耦合：好的顺序取决于载体将处于何处，而载体的位置又取决于顺序。两阶段方法通过在两个层级之间交替来处理这种耦合。我们的表可以看作一次这样的交替步骤，被冻结下来并在整个搜索中复用。}

\subsection{Deep reinforcement learning for routing\quad 用于路径问题的深度强化学习}

Most deep reinforcement learning methods for routing are developed for the travelling salesman problem (TSP) and the capacitated vehicle routing problem, in which trucks with limited load serve customers. In both problems, customers are fixed points and a route is an ordering of these points. Learning to construct routes began with pointer networks \citep{vinyals2015} and was extended to vehicle routing by \citet{nazari2018}. The attention model of \citet{kool2019} established the encoder--decoder architecture that most later work adopts. POMO \citep{kwon2020} exploits the multiple symmetric starts of a route, and Sym-NCO \citep{kim2022} and efficient active search \citep{hottung2022} add symmetry training and instance-level adaptation. Improvement-based policies \citep{wu2022} and heavy-decoder models \citep{luo2023} target better quality and larger sizes. Surveys by \citet{bengio2021} and \citet{mazyavkina2021} cover the field, and \citet{wuy2026} apply it to clustered vehicle routing in this journal. In these problems, the cost of a route is the sum of its edge lengths, and the load limit is checked by simple comparisons. A reward therefore costs one pass over the route, and no optimisation problem has to be solved.

\zh{路径问题的深度强化学习方法大多针对旅行商问题（TSP）和带容量约束的车辆路径问题（即载重有限的卡车服务客户）。两类问题中客户都是固定的点，一条路线就是这些点的一个排序。学习构造路线始于指针网络 \citep{vinyals2015}，\citet{nazari2018} 将其扩展到车辆路径问题。\citet{kool2019} 的注意力模型确立了此后大多数工作采用的编码器—解码器结构。POMO \citep{kwon2020} 利用路线的多个对称起点，Sym-NCO \citep{kim2022} 和高效主动搜索 \citep{hottung2022} 加入了对称性训练和实例级自适应。基于改进的策略 \citep{wu2022} 和重解码器模型 \citep{luo2023} 追求更好的质量和更大的规模。\citet{bengio2021} 和 \citet{mazyavkina2021} 的综述覆盖了这一领域，\citet{wuy2026} 在本刊将其应用于聚类车辆路径问题。在这些问题中，一条路线的成本就是各段边长之和，载重限制只需简单比较即可检查。因此计算一次奖励只需遍历一遍路线，不需要求解任何优化问题。}

The CVTSP combines two decisions: a visiting order, as in the TSP, and, for that order, the continuous points where the drone leaves and rejoins the carrier for each target. A visiting order alone therefore does not fix the cost of a route; the second decision requires solving an SOCP. DRL methods for truck--drone routing \citep{bogyrbayeva2023,peng2025} restrict launch and landing to network nodes, so their reward is still computed directly. Fewer learning methods address problems with continuous launch and landing points. Two studies learn policies for mothership problems, in which drones are launched from and recovered by a moving carrier. \citet{yan2026} study a truck that serves as a moving launch and recovery platform for two drones delivering to moving targets, and train a multi-agent proximal policy optimisation method. \citet{liu2026} apply proximal policy optimisation to eVTOL-and-drone routing with no-fly zones; the eVTOL chooses among a finite set of moves on a grid, and the policy is retrained for each instance. The single-drone mothership problem is the CVTSP itself \citep{erdogan2021}, so we do not claim that learning methods are absent from this class of problems. In both studies, however, the reward is computed by the simulated environment from elapsed time, distance and penalty terms, without solving an optimisation problem; the take-off and landing points result from the policy's step-by-step moves in the simulation. In the CVTSP formulation used here, the policy chooses only the visiting order, and the take-off and landing points are then optimised exactly for that order by solving an SOCP, so the cost of an order is known only after this solve. Our argument concerns the cost of training when every reward requires such a solve.

\zh{CVTSP 包含两层决策：一是访问顺序，与 TSP 相同；二是在给定顺序下，无人机对每个目标离开和返回载体的连续位置。因此仅凭访问顺序无法确定路线成本，第二层决策需要求解一个 SOCP。卡车—无人机路径问题的 DRL 方法 \citep{bogyrbayeva2023,peng2025} 把起飞和降落限制在网络节点上，因此其奖励仍可直接计算。针对起降点连续的问题，学习方法较少。有两项研究为母舰类问题学习策略，即无人机从移动载体起飞并由其回收。\citet{yan2026} 研究一辆卡车作为两架无人机的移动起降平台、向移动目标配送的问题，并训练一种多智能体近端策略优化方法。\citet{liu2026} 将近端策略优化用于带禁飞区的 eVTOL 与无人机路径问题；eVTOL 在网格上从有限的移动动作中选择，且策略需针对每个实例重新训练。单无人机的母舰问题就是 CVTSP 本身 \citep{erdogan2021}，因此我们不主张这类问题中没有学习方法。但在这两项研究中，奖励都由仿真环境根据经过的时间、距离和惩罚项计算，不需要求解优化问题；起飞点和降落点由策略在仿真中逐步的移动动作产生。而在本文采用的 CVTSP 形式中，策略只选择访问顺序，起降点随后通过求解一个 SOCP 针对该顺序精确优化，因此一个顺序的成本只有在这次求解之后才能知道。我们的论点关注的是每个奖励都需要这样一次求解时的训练成本。}

\subsection{Expensive objectives and self-improvement\quad 昂贵目标与自我改进}

In some combinatorial problems, evaluating the objective of a single solution is itself costly, because it requires a simulation, an experiment or an optimisation solve; in the CVTSP it requires an SOCP. Learning methods for such problems either reduce the number of evaluations or replace evaluations with a cheaper estimate. \citet{kim2024} take the first approach: symmetric replay training reuses solutions that have already been evaluated through their symmetries, and thereby saves evaluations. It offers no estimate for a route that was never evaluated. The second approach learns a surrogate. \citet{sun2026} train a neural network offline to estimate drone flight times from trajectory-optimisation data and use it inside an order-first, split-second heuristic for truck and drone delivery. SurCo \citep{ferber2023} learns linear surrogate costs for nonlinear combinatorial problems. These methods learn or reuse a cost signal, and none of those we read provides a guarantee on routes outside the evaluated set. A learned surrogate can be accurate on average, but it can also be wrong in a direction that is hard to detect, and its training set has to be produced in advance. A bound derived from the solution of the instance at hand needs no training data and has a known direction of error, which makes it usable as a screening rule, although in our runs it accepts relatively few routes on its own (Section~\ref{sec:search}).

\zh{在一些组合优化问题中，计算单个解的目标值本身就很昂贵，因为它需要一次仿真、实验或优化求解；在 CVTSP 中需要求解一个 SOCP。针对这类问题的学习方法要么减少评估次数，要么用更便宜的估计代替评估。\citet{kim2024} 采用第一种思路：对称回放训练利用对称性重用已评估过的解，从而节省评估次数。它不能对从未评估过的路线给出估计。第二种思路是学习代理模型。\citet{sun2026} 用轨迹优化数据离线训练一个神经网络来估计无人机飞行时间，并在卡车—无人机配送的“先排序、后拆分”启发式中使用它。SurCo \citep{ferber2023} 为非线性组合问题学习线性代理代价。这些方法学习或复用代价信号，而在我们读过的方法中，没有一个能对已评估集合之外的路线提供保证。学习得到的代理在平均意义上可能准确，但也可能在难以察觉的方向上出错，并且其训练集必须事先生成。由当前实例的解导出的上界无需训练数据，误差方向已知，因此可以用作筛选规则，尽管在我们的实验中它单独接受的路线相对较少（第~\ref{sec:search} 节）。}

Self-improvement offers a second route to lower training cost, in which a policy trains on labels that it produces itself. \citet{luo2023} train by supervised learning on optimal solutions of small instances; \citet{luo2024} remove the need for such labels by letting the model improve its own solutions, and \citet{pirnay2024} sample several solutions from the current model and imitate the best one. \citet{luttmann2025} extend the idea to problems in which several agents, such as vehicles or machines, act together. On the problems these papers study, comparing two candidate solutions is cheap, so labels can be accepted or rejected by direct evaluation. In the CVTSP, the same test costs one SOCP per candidate. Our method keeps the self-improvement loop and uses the bound to decide which candidates deserve that cost. This changes the economics of the loop. Generating a hundred candidates is cheap, but verifying a hundred candidates is not, so the screening step becomes the part of the method that determines its budget.

\zh{自我改进是降低训练成本的另一条途径，即策略用自己产生的标签进行训练。\citet{luo2023} 在小规模实例的最优解上进行监督学习；\citet{luo2024} 让模型改进自己的解，从而不再需要这类标签；\citet{pirnay2024} 从当前模型采样多个解并模仿其中最好的一个。\citet{luttmann2025} 将这一思路扩展到多个智能体（如多辆车或多台机器）协同行动的问题。在这些论文研究的问题上，比较两个候选解很廉价，因此可以直接评估来接受或拒绝标签。在 CVTSP 中，同样的检验需要为每个候选求解一次 SOCP。我们的方法保留自我改进循环，并用上界决定哪些候选值得付出这一代价。这改变了循环的经济性：生成一百个候选很廉价，验证一百个候选却不然，因此筛选步骤成为决定方法预算的部分。}

\subsection{Research gap\quad 研究空白}
\label{sec:gap}

Table~\ref{tab:gap} compares the closest work on five properties, which are defined in the table note. Among the sources we read, none combines a continuous route-cost subproblem, a training signal for routes that were never solved, and a primal upper bound behind that signal. The table shows two separate groups: optimisation methods for the CVTSP and the mothership problem handle the continuous subproblem but learn no policy, whereas learning methods for routing and drone routing build routes with a policy but avoid the continuous subproblem or replace it by simulation. Two parts of our design have direct precedents in the CVTSP literature. Evaluating other orders with the sortie points of a solved solution appears in the local search of \citet{erdogan2021}, and ranking orders by a cheap bound before solving them appears in the exact method of \citet{gambella2018}. The contrast with \citet{li2025} concerns the matrix. Their approximate matrix is initialised in closed form and, after a fixed-order problem is solved, updated only on traversed arcs; the other entries keep their prior values. The matrix guides a heuristic and is regenerated for diversification. Our table is built from one primal solution in a single step, assigns a valid value to every arc and bounds the cost of every order. Given these precedents, we claim as new neither the upper bound itself (Proposition~\ref{prop:bound} in Section~\ref{sec:bound}), which follows from keeping the take-off and landing points of one solved order, nor the idea of assigning a cost to an order without solving its SOCP. Our claim concerns their use to train a DRL policy: as a reward for routes that are never solved and as a ranking of candidate imitation labels, each of which is verified exactly before it is used.

\zh{表~\ref{tab:gap} 从五个性质比较了最接近的工作，这些性质在表注中定义。在我们读过的文献中，没有一项同时具备连续的路线代价子问题、针对从未求解路线的训练信号，以及支撑该信号的原始上界。表中可以看到两个彼此分开的群体：CVTSP 和母舰问题的优化方法处理连续子问题，但不学习策略；路径问题和无人机路径问题的学习方法用策略构造路线，但回避了连续子问题或以仿真代替它。我们设计中的两部分在 CVTSP 文献中有直接先例：用已求解解的架次点评估其他顺序见于 \citet{erdogan2021} 的局部搜索，在求解前用廉价界对顺序排序见于 \citet{gambella2018} 的精确方法。与 \citet{li2025} 的区别在于矩阵。他们的近似矩阵以闭式初始化，在求解一个定序问题后只更新经过的弧，其余元素保持先前的值。该矩阵用于引导启发式搜索，并为多样化而重新生成。我们的表由一个原始解一步构建，为每条弧赋予有效值，并为每个顺序的代价提供上界。鉴于这些先例，我们不主张以下两点是新的：一是上界本身（第~\ref{sec:bound} 节命题~\ref{prop:bound}：保留一个已求解顺序的起飞点和降落点，就能得到任意顺序代价的上界）；二是不解 SOCP、直接给一个顺序算出代价的做法。我们的主张在于把二者用于训练 DRL 策略：作为从未求解路线的奖励，以及作为候选模仿标签的排序，每个标签在使用前都经过精确验证。}

\begin{table}[!htbp]
\centering
\caption{Closest work compared on properties relevant to this paper. Columns: (a) the route cost is the optimal value of a continuous optimisation problem (``nodes'': launch and landing restricted to network nodes; ``simulated'': cost accumulated by a simulator); (b) the signal a learning method trains on, or n/a if nothing is learned; (c) how a route whose subproblem is not solved receives a value; (d) a primal upper bound on the cost of every order, obtained from one solved solution; (e) a learned policy constructs the route.\protect\newline\zhcap{与本文相关性质上的最接近工作比较。各列：(a) 路线代价是一个连续优化问题的最优值（“nodes”：起降限制在网络节点上；“simulated”：代价由仿真器累积）；(b) 学习方法所用的训练信号，若不学习则为 n/a；(c) 子问题未被求解的路线如何获得取值；(d) 由一个已求解解得到的、对每个顺序代价的原始上界；(e) 由学习得到的策略构造路线。}}
\label{tab:gap}
\small
\resizebox{\textwidth}{!}{\begin{tabular}{lccccc}
\toprule
Work & \begin{tabular}[c]{@{}c@{}}(a) Continuous\\ subproblem\end{tabular} & \begin{tabular}[c]{@{}c@{}}(b) Training\\ signal\end{tabular} & \begin{tabular}[c]{@{}c@{}}(c) Value for\\ unsolved routes\end{tabular} & \begin{tabular}[c]{@{}c@{}}(d) Primal bound\\ on every order\end{tabular} & \begin{tabular}[c]{@{}c@{}}(e) Learned policy\\ builds the route\end{tabular} \\
\midrule
\multicolumn{6}{l}{\textit{CVTSP and mothership optimisation}} \\
\citet{gambella2018} & Yes & n/a & Lower bound & No & No \\
\citet{poikonen2020} & Yes & n/a & Lower bound & No & No \\
\citet{erdogan2021}, ILS & Yes & n/a & Retained sortie points & Neighbours only & No \\
\citet{li2025} & Yes & n/a & Estimate & No & No \\
\citet{vieira2026} & Yes & n/a & Lower bound & No & No \\
\midrule
\multicolumn{6}{l}{\textit{Learning for routing and drone routing}} \\
\begin{tabular}[l]{@{}l@{}}\citet{kool2019};\\ \citet{kwon2020}\end{tabular} & No & Exact objective & Closed form & n/a & Yes \\
\begin{tabular}[l]{@{}l@{}}\citet{bogyrbayeva2023};\\ \citet{peng2025}\end{tabular} & No (nodes) & Exact objective & Closed form & n/a & Yes \\
\citet{liu2026} & No (simulated) & Simulated reward & Simulated & n/a & Yes \\
\midrule
\multicolumn{6}{l}{\textit{Learning under expensive evaluation}} \\
\citet{kim2024} & No & Exact, with replay & Not estimated & n/a & Yes \\
\begin{tabular}[l]{@{}l@{}}\citet{pirnay2024};\\ \citet{luo2024};\\ \citet{luttmann2025}\end{tabular} & No & Exact objective & Closed form & n/a & Yes \\
\midrule
This paper & Yes & Exact solves + bound & Retained sortie points & Yes & Yes \\
\bottomrule
\end{tabular}}
\end{table}

\section{Problem formulation\quad 问题描述}
\label{sec:problem}

Let $d \in \mathbb{R}^{2}$ be the depot and $p_{1},\dots,p_{J} \in \mathbb{R}^{2}$ the targets. The carrier moves at speed $v_{c}$ and the drone at speed $v_{d} > v_{c}$. The drone can fly for at most $E$ time units per sortie. Each target must be visited in exactly one sortie. For target $i$, the drone leaves the carrier at take-off point $s_{i}$, flies to $p_{i}$ and returns to the carrier at landing point $l_{i}$. The carrier travels from $s_{i}$ to $l_{i}$ during the sortie. The depot is both origin and destination, as in all benchmark instances we use.

\zh{设 $d \in \mathbb{R}^{2}$ 为基地，$p_{1},\dots,p_{J} \in \mathbb{R}^{2}$ 为目标。载体速度为 $v_{c}$，无人机速度为 $v_{d} > v_{c}$。无人机每个架次最多飞行 $E$ 个时间单位。每个目标必须恰好在一个架次中被访问。对于目标 $i$，无人机在起飞点 $s_{i}$ 离开载体，飞到 $p_{i}$，并在降落点 $l_{i}$ 回到载体；架次期间载体从 $s_{i}$ 行驶到 $l_{i}$。与我们使用的所有基准实例一样，基地既是起点也是终点。}

\subsection{The fixed-order problem\quad 定序问题}

Let $\pi = (\pi_{1},\dots,\pi_{J})$ be a visiting order, and set $\pi_{0} = \pi_{J+1} = 0$ with $l_{0} = s_{0} = d$, so that the first leg starts and the last leg ends at the depot. The decision variables are the take-off and landing points $s_{i}, l_{i} \in \mathbb{R}^{2}$, the outbound and inbound flight times $t^{1}_{i}$ and $t^{2}_{i}$, the sortie duration $\tau_{i}$, and the carrier travel time $T_{k}$ on leg $k$. Sortie variables are indexed by target, and leg variables by position in the order. The carrier--vehicle problem for order $\pi$ is
\begin{align}
f(\pi) = \min \quad & \sum_{i=1}^{J} \tau_{i} + \sum_{k=0}^{J} T_{k} \label{eq:obj} \\
\text{s.t.} \quad & \norm{s_{i} - p_{i}} \le v_{d}\, t^{1}_{i}, && i = 1,\dots,J, \label{eq:out}\\
& \norm{l_{i} - p_{i}} \le v_{d}\, t^{2}_{i}, && i = 1,\dots,J, \label{eq:in}\\
& \norm{l_{i} - s_{i}} \le v_{c}\, \tau_{i}, && i = 1,\dots,J, \label{eq:sync}\\
& \tau_{i} = t^{1}_{i} + t^{2}_{i}, \quad \tau_{i} \le E, && i = 1,\dots,J, \label{eq:end}\\
& \norm{s_{\pi_{k+1}} - l_{\pi_{k}}} \le v_{c}\, T_{k}, && k = 0,\dots,J, \label{eq:leg}\\
& t^{1}_{i},\, t^{2}_{i},\, \tau_{i},\, T_{k} \ge 0 . \label{eq:nonneg}
\end{align}
Every constraint is a second-order cone or a linear constraint, so the problem is an SOCP \citep{lobo1998}. Constraints \eqref{eq:out} and \eqref{eq:in} bound the flight times from below by the flight distances, so the drone may fly more slowly or wait in the air; $\tau_{i}$ is therefore the time between take-off and landing, and \eqref{eq:sync} requires the carrier to reach $l_{i}$ from $s_{i}$ within that time. At an optimal solution, $\tau_{i}$ equals the larger of the drone's and the carrier's travel times in the sortie, as in \citet{erdogan2021}. The CVTSP asks for $\min_{\pi} f(\pi)$. Mission completion time is the objective in \eqref{eq:obj}, because the carrier is always either on a leg or within a sortie. Table~\ref{tab:notation} lists the main notation.

\zh{设 $\pi = (\pi_{1},\dots,\pi_{J})$ 为一个访问顺序，令 $\pi_{0} = \pi_{J+1} = 0$，且 $l_{0} = s_{0} = d$，即第一段从基地出发、最后一段回到基地。决策变量为起飞点和降落点 $s_{i}, l_{i} \in \mathbb{R}^{2}$、去程和回程飞行时间 $t^{1}_{i}$ 与 $t^{2}_{i}$、架次持续时间 $\tau_{i}$，以及第 $k$ 段的载体行驶时间 $T_{k}$。架次变量按目标编号，行驶段变量按在顺序中的位置编号。顺序 $\pi$ 的载体—无人机问题如上式所示。每个约束都是二阶锥约束或线性约束，因此该问题是一个 SOCP \citep{lobo1998}。约束 \eqref{eq:out} 和 \eqref{eq:in} 只要求飞行时间不小于飞行距离所需的时间，因此无人机可以减速飞行或在空中等待；$\tau_{i}$ 就是从起飞到降落的时间，约束 \eqref{eq:sync} 要求载体在这段时间内从 $s_{i}$ 行驶到 $l_{i}$。在最优解中，$\tau_{i}$ 等于该架次中无人机与载体两者行程时间中较大的那个，这与 \citet{erdogan2021} 一致。CVTSP 求 $\min_{\pi} f(\pi)$。式~\eqref{eq:obj} 的目标就是任务完成时间，因为载体总是处于某个航段上或某个架次中。表~\ref{tab:notation} 列出了主要符号。}

\begin{table}[!htbp]
\centering
\caption{Main notation.\protect\newline\zhcap{主要符号。}}
\label{tab:notation}
\small
\begin{tabular}{ll}
\toprule
Symbol & Meaning \\
\midrule
$J$ & Number of targets \\
$d$, $p_{i}$ & Depot and target locations \\
$v_{c}$, $v_{d}$ & Carrier and drone speeds \\
$E$ & Maximum sortie duration \\
$\pi$ & Visiting order, with $\pi_{0} = \pi_{J+1} = 0$ denoting the depot \\
$s_{i}$, $l_{i}$ & Take-off and landing points of the sortie for target $i$ \\
$t^{1}_{i}$, $t^{2}_{i}$, $\tau_{i}$ & Outbound time, inbound time and duration of a sortie \\
$T_{k}$ & Carrier travel time on leg $k$ \\
$f(\pi)$ & Optimal value of the fixed-order problem for $\pi$ \\
$q$ & Anchor route, solved exactly \\
$C^{q}$, $F^{q}$ & Directed table and total sortie time from the solution for $q$ \\
$U_{q}(\pi)$ & Certified upper bound on $f(\pi)$ \\
\bottomrule
\end{tabular}
\end{table}

\subsection{Cost structure\quad 代价结构}

Two features of this formulation drive the method. First, constraints \eqref{eq:out}--\eqref{eq:end} involve the variables of a single target and make no reference to $\pi$. Second, the order appears only in \eqref{eq:leg}, which links the landing point of one target to the take-off point of the next. Any solution that fixes all sortie variables therefore induces a closed-form cost for every order, since the leg times are the only remaining quantities. Evaluating a fixed-order problem with $J$ targets takes tens of milliseconds at moderate sizes in our implementation, about 0.03\,s per call at $J \approx 60$, while evaluating a table sum takes microseconds. The number of orders that can be scored in one solve time differs by several orders of magnitude, which is the leverage the method relies on. This is the structure that the next section exploits. The policy in Section~\ref{sec:method} selects $\pi$, and the SOCP determines the sortie points, so the method follows the natural decomposition into a discrete upper level and a continuous lower level. Only the upper level is learned.

\zh{该模型的两个特征驱动了我们的方法。第一，约束 \eqref{eq:out}--\eqref{eq:end} 只涉及单个目标的变量，与 $\pi$ 无关。第二，顺序只出现在 \eqref{eq:leg} 中，它把一个目标的降落点与下一个目标的起飞点联系起来。因此，任何固定了全部架次变量的解都为每个顺序诱导出一个闭式代价，因为剩下的只有航段时间。在我们的实现中，求解一个有 $J$ 个目标的定序问题在中等规模下需要数十毫秒，$J \approx 60$ 时约 0.03\,s 一次，而计算表求和只需几微秒。在一次求解的时间内能打分的顺序数相差几个数量级，这正是该方法依赖的杠杆。下一节利用的就是这一结构。第~\ref{sec:method} 节中的策略选择 $\pi$，SOCP 确定架次点，因此方法遵循离散上层与连续下层的自然分解，只有上层是学习得到的。}

\section{Methodology\quad 方法}
\label{sec:method}

\subsection{Overview\quad 概览}

Figure~\ref{fig:method} shows one training step, and Table~\ref{alg:train} lists it step by step for a batch of instances. A frozen reference policy samples 200 candidate routes per instance. A small edge predictor selects one of them as the anchor, and the anchor is solved exactly. The solution yields a table $C$ that scores every candidate. Local search on the table produces an expert route, which is verified exactly before it can serve as an imitation label. The learning policy then receives two updates: a reinforcement step on rewards drawn mainly from the table, and an imitation step on the verified experts. The number of exact solves per instance is independent of the number of POMO routes, which is what keeps the solver budget small. In POMO with an exact reward, a batch of $B$ instances with $P$ routes costs $BP$ solves. Here it costs at most $3B$ (one anchor, one policy route and at most one expert check per instance), plus at most $2B$ in refresh rounds for instances with $J \ge 60$, regardless of $P$; the training run averaged 2.89 solves per instance.

\zh{图~\ref{fig:method} 展示了一个训练步骤，表~\ref{alg:train} 针对一批实例逐步列出了该步骤。冻结的参考策略为每个实例采样 200 条候选路线。一个小型边预测器从中选出一条作为锚点，并对锚点做精确求解。该解产生一张表 $C$，用来给每个候选打分。在表上的局部搜索产生一条专家路线，它在用作模仿标签之前要经过精确验证。随后学习策略接受两种更新：基于主要来自表的奖励的强化学习步骤，以及基于已验证专家的模仿步骤。每个实例的精确求解次数与 POMO 路线数无关，这正是求解器预算保持较小的原因。在使用精确奖励的 POMO 中，一批 $B$ 个实例、每个 $P$ 条路线需要 $BP$ 次求解；这里最多 $3B$ 次（每个实例一次锚点、一次策略路线、至多一次专家验证），对 $J \ge 60$ 的实例再加至多 $2B$ 次刷新，与 $P$ 无关；训练实际平均每个实例 2.89 次。}

\input{fig_method_bi}

\begin{table}[!htbp]
\caption{One training step for a batch of instances. Exact solves are marked.\protect\newline\zhcap{一批实例的一个训练步骤。精确求解已标出。}}
\label{alg:train}
\small
\hrule\smallskip
\noindent For each instance in the batch:
\begin{enumerate}
\setlength{\itemsep}{0pt}
\item Sample 200 routes from the frozen reference policy. \zhitem{从冻结的参考策略采样 200 条路线。}
\item Rank them with the edge predictor and choose the anchor $q$. \zhitem{用边预测器对其排序并选出锚点 $q$。}
\item \textbf{Exact solve:} solve the SOCP for $q$ and build the table $C^{q}$. \zhitem{\textbf{精确求解：}求解 $q$ 的 SOCP 并构建表 $C^{q}$。}
\item Score all sampled routes with $U_{q}$ and select 30 starts. \zhitem{用 $U_{q}$ 为所有采样路线打分，并选出 30 个起点。}
\item Run 2-opt on $C^{q}$ from each start, and LKH from a share of the starts when $J \ge 60$. \zhitem{在 $C^{q}$ 上从每个起点运行 2-opt；当 $J \ge 60$ 时，部分起点还运行 LKH。}
\item \textbf{Exact solve:} if the best route has $U_{q}$ below $f(q)$, solve it. Otherwise, with the near-miss probability, solve the closest route. \zhitem{\textbf{精确求解：}若最佳路线的 $U_{q}$ 低于 $f(q)$，则求解它；否则以近失概率求解最接近阈值的路线。}
\item \textbf{Exact solves:} for $J \ge 60$ after epoch 60, run up to two refresh rounds of the search, with one solve per round. \zhitem{\textbf{精确求解：}对于 $J \ge 60$ 且在第 60 轮之后，最多运行两轮刷新搜索，每轮求解一次。}
\end{enumerate}
\noindent For the batch as a whole:
\begin{enumerate}
\setlength{\itemsep}{0pt}
\setcounter{enumi}{7}
\item Decode the POMO routes of the learning policy and score them with the tables. \zhitem{解码学习策略的 POMO 路线，并用表为其打分。}
\item \textbf{Exact solve:} solve the best-scoring route of each instance. \zhitem{\textbf{精确求解：}求解每个实例中得分最好的路线。}
\item Update the policy by REINFORCE on the hybrid reward. \zhitem{用 REINFORCE 基于混合奖励更新策略。}
\item Update the policy by cross-entropy on the verified expert routes. \zhitem{用交叉熵基于已验证的专家路线更新策略。}
\item Update the edge predictor on the tables. \zhitem{用表更新边预测器。}
\end{enumerate}
\hrule
\end{table}

\subsection{A certified upper bound on all visiting orders\quad 对所有访问顺序的可证上界}
\label{sec:bound}

Let $q$ be a visiting order and let $(s^{*},l^{*},\tau^{*},T^{*})$ be an optimal solution of the fixed-order problem for $q$, with each sortie variable indexed by target. Write $F^{q} = \sum_{i} \tau^{*}_{i}$ and define a directed table on the nodes $\{0,1,\dots,J\}$, with node $0$ the depot:
\begin{equation}
C^{q}_{0j} = \frac{\norm{s^{*}_{j} - d}}{v_{c}}, \qquad
C^{q}_{i0} = \frac{\norm{d - l^{*}_{i}}}{v_{c}}, \qquad
C^{q}_{ij} = \frac{\norm{s^{*}_{j} - l^{*}_{i}}}{v_{c}} .
\end{equation}
For any order $\pi$, define the route value
\begin{equation}
U_{q}(\pi) = F^{q} + \sum_{k=0}^{J} C^{q}_{\pi_{k}\pi_{k+1}}, \qquad \pi_{0} = \pi_{J+1} = 0 .
\end{equation}

\zh{设 $q$ 为一个访问顺序，$(s^{*},l^{*},\tau^{*},T^{*})$ 为 $q$ 的定序问题的一个最优解，每个架次变量按目标编号。记 $F^{q} = \sum_{i} \tau^{*}_{i}$，并在节点 $\{0,1,\dots,J\}$（节点 0 为基地）上定义一张有向表，如上式所示。对任意顺序 $\pi$，路线值 $U_{q}(\pi)$ 的定义亦如上式。}

\begin{proposition}
\label{prop:bound}
For every visiting order $\pi$, $f(\pi) \le U_{q}(\pi)$. Moreover, $f(q) = U_{q}(q)$.
\end{proposition}

\zh{\textbf{命题~\ref{prop:bound}}\quad 对每个访问顺序 $\pi$，有 $f(\pi) \le U_{q}(\pi)$。此外，$f(q) = U_{q}(q)$。}


\begin{proof}
Fix $\pi$. Keep $s^{*}_{i}$, $l^{*}_{i}$, $t^{1*}_{i}$, $t^{2*}_{i}$ and $\tau^{*}_{i}$ for each target and set $T_{k} = C^{q}_{\pi_{k}\pi_{k+1}}$. Constraints \eqref{eq:out}--\eqref{eq:end} involve only these sortie variables and do not depend on the order, so they hold because they hold for $q$. Constraint \eqref{eq:leg} holds with equality by construction, and non-negativity is immediate. The point is thus feasible for the fixed-order problem of $\pi$, with objective value $U_{q}(\pi)$. Since $f(\pi)$ is the minimum over feasible points, $f(\pi) \le U_{q}(\pi)$. For $\pi = q$, each $T^{*}_{k}$ is at its lower bound at an optimal solution because the objective strictly increases in $T_{k}$, which gives $U_{q}(q) = f(q)$.
\end{proof}

\zh{\textbf{证明.}\quad 固定 $\pi$。对每个目标保留 $s^{*}_{i}$、$l^{*}_{i}$、$t^{1*}_{i}$、$t^{2*}_{i}$ 和 $\tau^{*}_{i}$，并令 $T_{k} = C^{q}_{\pi_{k}\pi_{k+1}}$。约束 \eqref{eq:out}--\eqref{eq:end} 只涉及这些架次变量且与顺序无关，它们对 $q$ 成立，因此仍然成立。约束 \eqref{eq:leg} 按构造以等号成立，非负性显然。因此该点对 $\pi$ 的定序问题可行，目标值为 $U_{q}(\pi)$。由于 $f(\pi)$ 是可行点上的最小值，$f(\pi) \le U_{q}(\pi)$。对于 $\pi = q$，由于目标函数关于 $T_{k}$ 严格递增，最优解中每个 $T^{*}_{k}$ 都处于其下界，从而 $U_{q}(q) = f(q)$。}


The proof does not use the fact that the route ends at the depot. If the destination differs from the origin, as the general problem allows \citep{erdogan2021}, the proposition holds with the arcs into node~$0$ measured to the destination. Three conditions apply in practice. The anchor must be solved to optimality; a merely feasible solution would still give a valid upper bound, but the equality at $q$ could fail. In the implementation, every anchor solution passes an audit of all constraints with an absolute tolerance of $10^{-5}$ before its table is used. A small violation in one sortie can only be repaired by lengthening times, so the bound is valid up to solver feasibility tolerances accumulated over the $J$ sorties, not up to $10^{-5}$ itself. The training code also checks the inequality, with a tolerance of $2 \times 10^{-5}$, whenever an expert, near-miss or refresh route is solved exactly, and it stops on a violation, which would indicate a numerical or implementation fault. The final run completed without a stop.

\zh{证明没有用到路线回到基地这一点。如果终点与起点不同（一般问题允许这种情况 \citep{erdogan2021}），只需把进入节点~$0$ 的弧改为到终点的行驶时间，命题仍然成立。实践中有三个条件。锚点必须求解到最优；仅可行的解仍会给出有效上界，但在 $q$ 处的等式可能不成立。在实现中，每个锚点解在其表被使用之前，都要以 $10^{-5}$ 的绝对容差审核全部约束。一个架次中的微小违反只能通过延长时间来修复，因此上界在 $J$ 个架次累积的求解器可行性容差范围内有效，而不是在 $10^{-5}$ 本身的范围内。当专家、近失或刷新路线被精确求解时，训练代码还会以 $2 \times 10^{-5}$ 的容差检查该不等式，若违反则停止运行，这将表明存在数值或实现错误。最终的训练运行没有发生停止。}

\begin{remark}
\label{rem:combine}
When several orders of an instance have been solved, each solution gives its own table, and the minimum of their route values, $\min_{q} U_{q}(\pi)$, is again an upper bound on $f(\pi)$. The training step takes this minimum over the anchor and expert tables (Section~\ref{sec:learning}). An entry-wise minimum of the tables is not valid: it splices together sortie points from different solutions, so that the take-off and landing points of one target could differ between its incoming and outgoing arcs.
\end{remark}

\zh{\textbf{注~\ref{rem:combine}}\quad 一个实例解过几个顺序时，每个解都给出自己的表，它们路线值的最小值 $\min_{q} U_{q}(\pi)$ 仍是 $f(\pi)$ 的上界。训练步骤对锚点表和专家表取这个最小值（第~\ref{sec:learning}~节）。逐元素取各表的最小值不成立：它会把不同解的架次点拼在一起，使同一个目标进来的边和出去的边用的起降点不一样。}


When the depot is both origin and destination, an order and its reverse have the same optimal cost, $f(\overleftarrow{\pi}) = f(\pi)$, where $\overleftarrow{\pi}$ visits the targets of $\pi$ in reverse order: exchanging the take-off and landing points of every sortie, and its outbound and inbound times, turns a solution for one into a solution for the other. The table does not share this symmetry. Its entry $C^{q}_{ij}$ runs from the landing point of $i$ to the take-off point of $j$, so reversing a route changes its internal arcs, and $U_{q}(\overleftarrow{q})$ generally exceeds $f(q)$. Local search must therefore account for the reversal of internal arcs. A 2-opt move reverses a section of the route. In a symmetric problem only the two arcs at the ends of the section change, but here every arc inside the section is traversed in the opposite direction and changes its cost. Prefix sums of the forward and reverse arc costs along the current route give this change in constant time, so each 2-opt move is still scored in constant time, without any solver call.

\zh{当基地既是起点又是终点时，一个顺序和它的反向顺序最优成本相同，即 $f(\overleftarrow{\pi}) = f(\pi)$，其中 $\overleftarrow{\pi}$ 以相反顺序访问 $\pi$ 的目标：把每个架次的起飞点和降落点互换，出程和回程时间也互换，一个顺序的解就变成另一个顺序的解。表没有这种对称性。它的元素 $C^{q}_{ij}$ 是从 $i$ 的降落点到 $j$ 的起飞点，所以把路线反过来会改变内部的弧，$U_{q}(\overleftarrow{q})$ 一般大于 $f(q)$。因此局部搜索必须考虑内部弧的反转。一次 2-opt 移动会把路线的一段反过来。在对称问题中，只有这一段两端的两条弧会变；而这里段内每条弧都改为反向经过，成本也随之改变。沿当前路线预先计算正向和反向弧成本的前缀和，就能在常数时间内得到这一变化，所以每次 2-opt 移动仍可在常数时间内打分，无需任何求解器调用。}

\subsection{Policy, anchor selection and solver-free expert search\quad 策略、锚点选择与无需求解器的专家搜索}
\label{sec:search}

The learning policy is an attention encoder--decoder of the POMO type \citep{kwon2020}, with 4 encoder layers, embedding dimension 256, 16 attention heads and feed-forward dimension 512. Each node has six features: its coordinates relative to the depot, divided by the largest depot--target distance of the instance; its distance to the depot on the same scale; a depot indicator; the speed ratio $v_{c}/v_{d}$; and the drone reach $v_{d}E$ on the same scale. The speed and endurance features let one network serve instances with different operating parameters. The policy decodes one route per POMO start, and in training every target serves as a start, so $P = J$.

\zh{学习策略是 POMO 型的注意力编码器–解码器 \citep{kwon2020}，有 4 个编码层，嵌入维度 256，16 个注意力头，前馈维度 512。每个节点有六个特征：相对基地的坐标，除以该实例中基地到目标的最大距离；以同一尺度计的到基地距离；基地标记；速度比 $v_{c}/v_{d}$；以及以同一尺度计的无人机航程 $v_{d}E$。速度和续航特征使一个网络能够处理运行参数不同的实例。策略对每个 POMO 起点解码一条路线；训练中每个目标都作为起点，所以 $P = J$。}

Two policies interact during training. The learning policy is updated at every step. A frozen reference policy generates search material and is replaced only when the learning policy obtains a strictly better exact mean on a fixed validation set of 1,000 instances. This gate ensures that the expert search starts from routes of validated quality.

\zh{训练中有两个策略相互作用。学习策略在每一步都更新。冻结的参考策略生成搜索素材，只有当学习策略在一个固定的 1,000 个实例的验证集上取得严格更好的精确均值时才被替换。这一门控保证专家搜索从质量经过验证的路线出发。}

An edge predictor maps the detached encoder output to a non-negative directed table of size $(J+1) \times (J+1)$. It is trained with a Huber loss towards tables built from exact solutions, and its loss does not reach the policy. Its only role is to rank the 200 reference routes and choose the anchor, so a poor prediction costs search quality and cannot affect validity.

\zh{边预测器把切断梯度后的编码器输出映射为一张大小为 $(J+1) \times (J+1)$ 的非负有向表。它以 Huber 损失向由精确解构建的表回归，其损失不会传到策略。它唯一的作用是对 200 条参考路线排序并选出锚点，因此预测不准只会降低搜索质量，不会影响有效性。}

The anchor table scores all sampled routes. We select 30 starts: the 20 routes with the lowest bound, and 10 further routes chosen for diversity from a pool of the 60 next-best routes within a relative gap of 5\%. From each start, directed 2-opt runs on the table for at most 30 moves and accepts only strict decreases of $U_{q}$. Local search on the table is complementary to LKH \citep{lin1973,helsgaun2000}: on large instances, a share of the starts is also passed to LKH as an asymmetric TSP on $C^{q}$. In every case the search minimises a surrogate that can only overestimate the true cost. The LKH time budget for all participating starts is deterministic, $1.5 \times 10^{-4} J^{2}$ seconds scaled by the LKH share, so that the search cost of a batch can be predicted in advance.

\zh{锚点表为所有采样路线打分。我们选出 30 个起点：上界最低的 20 条路线，以及从相对差距在 5\% 以内的 60 条次优路线中按多样性选出的另外 10 条。从每个起点出发，有向 2-opt 在表上最多运行 30 步，只接受 $U_{q}$ 的严格下降。表上的局部搜索与 LKH \citep{lin1973,helsgaun2000} 互补：在大实例上，部分起点还会作为 $C^{q}$ 上的非对称 TSP 交给 LKH。无论哪种情况，搜索最小化的都是一个只会高估真实代价的代理。所有参与起点的 LKH 时间预算是确定性的，为 $1.5 \times 10^{-4} J^{2}$ 秒乘以 LKH 份额，因此一批的搜索成本可以事先预测。}

Proposition~\ref{prop:bound} gives the screening rule. If a route satisfies $U_{q}(\pi) < f(q)$, then $f(\pi) < f(q)$ holds without any solve, so the route is a certified improvement. Its exact cost is then computed once, and the route becomes the expert. If no route clears the threshold, the anchor remains the expert. Because the bound is conservative, some truly better routes fail this test. With a small probability, the near-miss rule therefore solves the route with the smallest relative gap above the threshold and accepts it if it is exactly better. The probability adapts to training progress: for small instances it moves between 1\% and 5\% as the imitation accuracy and the share of policy routes that beat the expert rise. Each instance has a budget of one such verification, and the solve is charged whether or not the route is accepted. If the exact cost of a near-miss route exceeds its bound, the run stops, because that would indicate a numerical or implementation fault.

\zh{命题~\ref{prop:bound} 给出筛选规则。若一条路线满足 $U_{q}(\pi) < f(q)$，则无需求解即有 $f(\pi) < f(q)$，该路线是一个可证的改进。随后计算一次其精确代价，该路线成为专家。若没有路线越过阈值，锚点仍为专家。由于上界是保守的，一些真正更好的路线会通不过这一检验。因此，近失规则以一个较小的概率求解高于阈值且相对差距最小的路线，若它精确地更好则接受。该概率随训练进展自适应：对小实例，它随模仿准确率以及策略路线击败专家的比例上升而在 1\% 到 5\% 之间变化。每个实例有一次这种验证的预算，无论路线是否被接受，这次求解都会计入。若近失路线的精确代价超过其上界，运行将停止，因为这将表明存在数值或实现错误。}

In practice the first-round test accepts relatively few experts once training has progressed, and most accepted experts come from near-miss checks and refresh rounds (Section~\ref{sec:recipe}). Through epoch 133, the first-round test led to 128,154 exact verifications, near-miss checks accepted 159,866 routes and refresh rounds accepted 162,888. At epoch 133, 659 of 10,000 instances passed the first-round test, whereas 3,788 received an improved expert. Near-miss routes checked at that epoch had bounds 14.6\% above the exact cost of the anchor on average, and 40.8\% of them were exactly better than the anchor. The bound therefore works less as a certificate than as a ranking that decides which few routes receive an exact solve. Its guarantee matters for validity: every accepted label is exact, whichever rule proposed it.

\zh{实践中，训练推进之后，第一轮检验单独接受的专家相对较少，大多数被接受的专家来自近失检验和刷新轮次（第~\ref{sec:recipe} 节）。截至第 133 轮，第一轮检验带来 128,154 次精确验证，近失检验接受了 159,866 条路线，刷新轮次接受了 162,888 条。在第 133 轮，10,000 个实例中有 659 个通过第一轮检验，而有 3,788 个得到了改进的专家。该轮被检验的近失路线，其上界平均比锚点的精确代价高 14.6\%，其中 40.8\% 精确地优于锚点。因此，上界与其说是一种证书，不如说是一种排序，决定哪几条路线值得一次精确求解。它的保证对有效性很重要：无论哪条规则提出了标签，每个被接受的标签都是精确的。}

\subsection{Learning from few exact solves\quad 基于少量精确求解的学习}
\label{sec:learning}

Each instance yields three kinds of cost information for the POMO routes $\pi^{1},\dots,\pi^{P}$ of the learning policy. The anchor table and the expert table bound every route, and the code takes the route-wise minimum of the two. The route with the lowest bound is then solved exactly. Writing $c^{p}$ for the exact cost of that route and for the bound of every other route, the hybrid advantage is
\begin{equation}
a^{p} = \frac{\bar{c} - c^{p}}{\max\{\sigma(c),\, 10^{-6}\}},
\end{equation}
where $\bar{c}$ and $\sigma(c)$ are the mean and standard deviation over the POMO routes of the instance. Standardising within each instance prevents hard instances from dominating the gradient. A second term compares the exactly solved policy route $\pi^{\mathrm{sel}}$ with the verified expert $\pi^{E}$,
\begin{equation}
a^{\mathrm{ex}} = \operatorname{clip}\!\left(\frac{f(\pi^{E}) - f(\pi^{\mathrm{sel}})}{|f(\pi^{E})|},\, -1,\, 1\right).
\end{equation}
The reinforcement loss combines the two terms using REINFORCE \citep{williams1992}:
\begin{equation}
\mathcal{L}_{\mathrm{RL}} = -\frac{1}{BP}\sum_{b,p} a^{p}_{b} \log p_{\theta}(\pi^{p}_{b}) - \frac{w}{B} \sum_{b} a^{\mathrm{ex}}_{b} \log p_{\theta}(\pi^{\mathrm{sel}}_{b}),
\end{equation}
where $B$ is the batch size. The weight $w$ rises linearly from 0 to 1 as the imitation accuracy of the previous epoch moves from 0.50 to 0.85, so the comparison with the expert begins only when the policy can produce competitive routes.

\zh{每个实例为学习策略的 POMO 路线 $\pi^{1},\dots,\pi^{P}$ 提供三类代价信息。锚点表和专家表都为每条路线给出上界，代码取二者中逐路线的较小值。上界最低的路线随后被精确求解。记 $c^{p}$ 为该路线的精确代价以及其他每条路线的上界，混合优势 $a^{p}$ 如上式所示，其中 $\bar{c}$ 和 $\sigma(c)$ 是该实例 POMO 路线上的均值和标准差。在每个实例内标准化可以防止困难实例主导梯度。第二项把被精确求解的策略路线 $\pi^{\mathrm{sel}}$ 与已验证的专家 $\pi^{E}$ 比较。强化学习损失用 REINFORCE \citep{williams1992} 组合这两项，其中 $B$ 为批大小。权重 $w$ 随上一轮的模仿准确率从 0.50 升至 0.85 而从 0 线性增至 1，因此与专家的比较只在策略能产生有竞争力的路线之后才开始。}

The hybrid advantage $a^{p}$ is biased, not only noisy. The solved route carries its exact cost while every other route carries a bound that exceeds its exact cost by a wide margin, so the solved route tends to receive the largest advantage by construction. In a diagnostic on 200 instances with two checkpoints of an earlier run of the method (epochs 100 and 365, not E133), the cosine between this advantage vector and the one obtained with all routes solved exactly was 0.61 and 0.67, falling to 0.53 and 0.54 at $J = 100$, and the route that was truly best received a negative advantage on 8\% and 12\% of the instances. Because the bound is tightest near the anchor, the bias can also favour routes that resemble the anchor; Section~\ref{sec:limits} returns to this point.

\zh{混合优势 $a^{p}$ 不仅有噪声，而且有偏。被求解的路线带着其精确代价，而其他每条路线带着的上界都大幅超过其精确代价，因此被求解的路线往往天然获得最大的优势。在一项针对 200 个实例、使用该方法早期一次运行的两个检查点（第 100 轮和第 365 轮，而非 E133）的诊断中，该优势向量与所有路线都精确求解时得到的优势向量之间的余弦相似度为 0.61 和 0.67，在 $J = 100$ 时降至 0.53 和 0.54，并且真正最好的路线在 8\% 和 12\% 的实例中得到了负优势。由于上界在锚点附近最紧，这种偏差也可能偏向与锚点相似的路线；第~\ref{sec:limits} 节将回到这一点。}

A separate optimiser step performs teacher-forced cross-entropy on the expert sequences. Each batch mixes current experts with a replay buffer of 50,000 experts, which supplies 30\% of the batch. Only verified experts enter this update, so an error in the bound can waste a solve but can never become a label. Imitation runs for five inner epochs at each update on batches of 1,000 sequences, split into micro-batches of 128. The two updates use a common network and take separate optimiser steps.

\zh{单独的一个优化器步骤对专家序列做教师强制的交叉熵训练。每批将当前专家与一个容量为 50,000 的专家回放缓冲区混合，后者占批次的 30\%。只有经过验证的专家进入这一更新，因此上界的误差最多浪费一次求解，永远不会变成标签。每次更新时，模仿学习在 1,000 条序列的批上运行五个内部轮次，并拆分为大小为 128 的微批。两种更新使用同一个网络，分别执行各自的优化器步骤。}

\subsection{Progressive search\quad 渐进式搜索}
\label{sec:recipe}

Search effort is scheduled by instance size and training stage. Instances with $J < 60$ keep the plain 2-opt search throughout. For $J \ge 60$, a ramp $r = \operatorname{clip}((e - 40)/60,\, 0,\, 1)$ at epoch $e$ raises the fraction of LKH starts to $r/2$ and raises the near-miss probability towards 1. After epoch 60, large instances also receive two refresh rounds: the most recently solved route supplies a new table, which is tight around that route, and the search continues on it. Each round solves at most one additional route.

\zh{搜索力度按实例规模和训练阶段调度。$J < 60$ 的实例始终使用普通 2-opt 搜索。对于 $J \ge 60$，第 $e$ 轮的斜坡 $r = \operatorname{clip}((e - 40)/60,\, 0,\, 1)$ 把 LKH 起点比例提高到 $r/2$，并把近失概率提高到接近 1。第 60 轮之后，大实例还会进行两轮刷新：最近求解的路线提供一张新表，它在该路线附近是紧的，搜索在其上继续。每轮最多额外求解一条路线。}

The quantity of imitation labels is normalised at the same time. An improved expert is repeated in the imitation batch so that the expected number of samples per instance stays near 1.44, with the repeat weight clipped to $[1, 5]$. Without this step, the imitation pressure would grow as the search improves.

\zh{与此同时，模仿标签的数量被归一化。改进后的专家在模仿批次中被重复，使每个实例的期望样本数保持在 1.44 左右，重复权重裁剪到 $[1, 5]$。没有这一步，模仿压力会随搜索改进而增长。}

The operator mix and size threshold were chosen in earlier development runs. The final recipe uses 2-opt and LKH without swap. The ramp over epochs 40 to 100 was chosen so that the policy first learns from the cheapest search and receives the stronger expert only once it has stabilised. We did not run a formal search over these values, and Section~\ref{sec:discussion} returns to this point.

\zh{算子组合和规模阈值在早期开发运行中选定。最终方案使用 2-opt 和 LKH，不使用 swap。第 40 至 100 轮的斜坡是这样选择的：策略先从最廉价的搜索中学习，在稳定之后才接受更强的专家。我们没有对这些取值做正式搜索，第~\ref{sec:discussion} 节会回到这一点。}

\section{Computational experiments\quad 计算实验}
\label{sec:experiments}

\subsection{Setup\quad 实验设置}

\paragraph{Baseline.（基线）} G200 is a POMO baseline that we previously trained for 200 epochs with an exact Gurobi reward on every sampled route, on 100,000 random instances per epoch drawn from the same instance generator as the proposed method. Its logged solver calls show that it decoded $J/4$ POMO starts per training instance during these 200 epochs, whereas the proposed method trains with $J$ starts; both models use the same start rule at test time. G200 uses a smaller network (6 encoder layers, embedding dimension 128, 8 heads, 1.27 million parameters, against 2.44 million for E133), a batch size of 512, and Adam with a constant learning rate of $5 \times 10^{-5}$ and weight decay $10^{-6}$. The two models therefore differ in architecture, training length and start count as well as in training signal, and the comparison in this paper is between complete training recipes. We use its saved epoch-200 checkpoint in this comparison. A separate retraining under the same baseline recipe on our workstation is reported in Sections~\ref{sec:efficiency} and~\ref{sec:ablation}.

\zh{\textbf{基线。}G200 是我们此前自行训练的 POMO 基线，对每条采样路线都使用精确 Gurobi 奖励，共训练 200 轮，每轮使用 100,000 个随机实例，来自与所提方法相同的实例生成器。其记录的求解器调用表明，在这 200 轮中它对每个训练实例解码 $J/4$ 个 POMO 起点，而所提方法训练时使用 $J$ 个起点；两个模型在测试时使用相同的起点规则。G200 使用更小的网络（6 个编码层、嵌入维度 128、8 个头、127 万参数，E133 为 244 万），批大小为 512，Adam 学习率恒为 $5 \times 10^{-5}$，权重衰减为 $10^{-6}$。因此，两个模型在网络结构、训练长度、起点数量以及训练信号上都不同，本文的比较是完整训练方案之间的比较。我们在本次比较中使用其第 200 轮保存的检查点。同一基线方案在本工作站上另行重训的结果见第~\ref{sec:efficiency} 节和第~\ref{sec:ablation} 节。}

\paragraph{Proposed model.（所提模型）} Training ran for 200 epochs of 10,000 instances each, with a batch size of 64, Adam with learning rate $10^{-4}$ and weight decay $10^{-6}$, and a cosine schedule with a 500-epoch horizon that was stopped at epoch 200. We report the checkpoint at epoch 133 (E133), selected by mean exact cost on a validation set of 2,000 instances, with 400 at each of $J \in \{20,40,60,80,100\}$. The selection used the original unaugmented validation protocol and no test data. One fixed-order SOCP is solved by Gurobi with one thread per call. Experiments ran on a workstation with an Intel i9-13900KF processor (32 hardware threads), 31\,GB of memory and an NVIDIA RTX 4090 GPU.

\zh{\textbf{所提模型。}训练共 200 轮，每轮 10,000 个实例，批大小 64，Adam 学习率 $10^{-4}$、权重衰减 $10^{-6}$，采用周期为 500 轮的余弦调度并在第 200 轮停止。我们报告第 133 轮的检查点（E133），它依据在一个 2,000 实例验证集上的平均精确代价选出，验证集在 $J \in \{20,40,60,80,100\}$ 各有 400 个实例。模型选择使用原来的无增强验证协议，没有使用测试数据。每个定序 SOCP 由 Gurobi 以单线程求解。实验在一台配有 Intel i9-13900KF 处理器（32 个硬件线程）、31\,GB 内存和 NVIDIA RTX 4090 GPU 的工作站上运行。}

\paragraph{Data.（数据）} Training instances have $J$ drawn uniformly from $\{20,\dots,100\}$ and targets drawn uniformly from the unit square. Speeds and endurance are sampled uniformly between the smallest and largest values in the literature benchmark, with $v_{c} = 0.2$, $v_{d} \in [0.4, 1]$ and $E \in [0.05, 1]$ in normalised units, so the speed ratio $v_{c}/v_{d}$ lies between 0.2 and 0.5. We evaluate on three sets. Real529 is the literature benchmark of \citet{gambella2018}, \citet{erdogan2021} and \citet{poikonen2020}, as compiled by \citet{li2025}, with 72, 132 and 325 instances. Its sizes include $J$ from 10 to 19 and $J = 200$, which lie outside the training range. It contains only two parameter settings, both at the edges of the training range: 325 instances with speed ratio 0.5 and endurance 20, and 204 with speed ratio 0.2 and endurance 1, in the units of the benchmark files. Test500 contains 100 new instances at each of $J \in \{20,40,60,80,100\}$. The generalisation set contains 100 new instances at each of $J \in \{150,200,300,500\}$. Test500 serves as the primary evidence because it matches the training distribution and allows a clean comparison at each size. We inspected test500 once, for the checkpoint E150 chosen by an earlier and coarser selection, before the refined selection that chose E133; the selection itself relied on the validation set only. Real529 is not a held-out set, because earlier checkpoints of this project were evaluated on it during method development.

\zh{\textbf{数据。}训练实例的 $J$ 在 $\{20,\dots,100\}$ 中均匀抽取，目标在单位正方形中均匀抽取。速度和续航在文献基准的最小值与最大值之间均匀采样，归一化单位下 $v_{c} = 0.2$、$v_{d} \in [0.4, 1]$、$E \in [0.05, 1]$，因此速度比 $v_{c}/v_{d}$ 介于 0.2 与 0.5 之间。我们在三个集合上评估。Real529 是 \citet{gambella2018}、\citet{erdogan2021} 和 \citet{poikonen2020} 的文献基准，由 \citet{li2025} 汇编，分别有 72、132 和 325 个实例。其规模包括训练范围之外的 $J = 10$ 至 19 以及 $J = 200$。它只包含两种参数设置，都位于训练范围的边缘：325 个实例速度比为 0.5、续航为 20，204 个实例速度比为 0.2、续航为 1（以基准文件的单位计）。Test500 在 $J \in \{20,40,60,80,100\}$ 各含 100 个新实例。泛化集在 $J \in \{150,200,300,500\}$ 各含 100 个新实例。Test500 是主要证据，因为它与训练分布一致，并允许在每个规模上做干净的比较。在细化后的选择选出 E133 之前，我们曾对较早、较粗的选择所选出的检查点 E150 查看过一次 test500；选择本身只依赖验证集。Real529 不是留出集，因为本项目的早期检查点在方法开发中曾在其上评估过。}

\paragraph{Test protocol.（测试协议）} All reported policy-quality comparisons use eight-view augmentation and greedy POMO, with $\min(J,100)$ starts per view: $J$ starts up to $J = 100$ and 100 starts on larger instances, including the generalisation set and the $J = 200$ instances of Real529. Each geometric view is decoded separately. Every distinct route is solved exactly and the best value over all views is reported. Reported policy costs are therefore exact for the decoded orders; the full mixed-integer solver's incumbents and lower bounds are reported separately. This reporting protocol does not change the original validation-based checkpoint selection.

\zh{\textbf{测试协议。}报告的所有策略解质量比较均使用八视角增强和 POMO 贪心解码，每个视角使用 $\min(J,100)$ 个起点：$J \le 100$ 时为 $J$ 个起点，更大的实例（包括泛化集和 Real529 中 $J = 200$ 的实例）为 100 个起点。每个几何视角单独解码。每条不同的路线都被精确求解，报告所有视角中的最好值。因此，报告的策略代价对解码出的顺序而言是精确的；完整混合整数求解器的当前最好可行解和下界单独报告。这一展示协议不改变原来的基于验证集的检查点选择。}

\paragraph{Metrics.（评价指标）} For model $A$ and baseline $B$ on instance $n$, the relative gap is $100\,(c^{A}_{n} - c^{B}_{n})/c^{B}_{n}$, so positive values mean that $A$ is worse. We report the mean over instances with a 95\% normal-approximation interval, together with the counts of wins, ties and losses at a tolerance of $10^{-3}$. We prefer the per-instance mean because a ratio of mean costs is dominated by the largest instances. On Real529, for example, the ratio of means is +0.12\%, whereas the per-instance mean is +0.21\%.

\zh{\textbf{评价指标。}对于模型 $A$ 与基线 $B$ 在实例 $n$ 上，相对差距为 $100\,(c^{A}_{n} - c^{B}_{n})/c^{B}_{n}$，正值表示 $A$ 更差。我们报告实例上的均值及 95\% 正态近似区间，以及在 $10^{-3}$ 容差下的胜、平、负计数。我们采用逐实例均值，因为平均代价之比受最大实例主导。例如在 Real529 上，均值之比为 +0.12\%，而逐实例均值为 +0.21\%。}

\subsection{Results on the literature benchmark\quad 文献基准上的结果}
\label{sec:real}

Table~\ref{tab:real} reports Real529 by size band. The mean gap to G200 is +0.21\%, with an interval that excludes zero. Gaps are below +0.53\% in every band up to $J = 100$. At $J = 200$, E133 is 1.25\% better than G200, with 21 wins out of 25. Ties are frequent at the smallest sizes, where 127 of 220 instances end at the same value, which suggests that both models often reach the same route. This holds for 25 instances from one source and lies outside the training range, so we treat it as an observation and not as evidence of general superiority: on the synthetic instances at the same size in Table~\ref{tab:gen}, E133 is behind.

\zh{表~\ref{tab:real} 按规模段报告 Real529。与 G200 的平均差距为 +0.21\%，区间不含零。在 $J = 100$ 及以下的每个规模段，差距都低于 +0.53\%。在 $J = 200$ 时，E133 比 G200 好 1.25\%，25 个实例中赢 21 个。最小规模上平局很常见，220 个实例中有 127 个得到相同的值，这说明两个模型经常得到相同的路线。$J = 200$ 的结果只来自一个来源的 25 个实例，且位于训练范围之外，因此我们把它视为一项观察，而不是普遍优势的证据：在表~\ref{tab:gen} 中同等规模的合成实例上，E133 是落后的。}

\begin{table}[!htbp]
\centering
\caption{Real529 with eight views, by size band. Gap is E133 relative to G200.\protect\newline\zhcap{八视角下按规模段的 Real529。差距为 E133 相对 G200。}}
\label{tab:real}
\small
\begin{tabular}{lrrrrrr}
\toprule
Band & $n$ & E133 & G200 & Gap (\%) & 95\% interval & Win/Tie/Loss \\
\midrule
$J = 10$--19 & 220 & 255.570 & 255.086 & +0.19 & [+0.12, +0.26] & 23/127/70 \\
$J = 20$--50 & 234 & 342.841 & 341.847 & +0.33 & [+0.21, +0.45] & 60/33/141 \\
$J = 100$ & 50 & 431.749 & 429.727 & +0.52 & [+0.13, +0.92] & 19/0/31 \\
$J = 200$ & 25 & 720.934 & 730.123 & $-$1.25 & [$-$1.78, $-$0.71] & 21/0/4 \\
All & 529 & 332.819 & 332.421 & +0.21 & [+0.13, +0.29] & 123/160/246 \\
\bottomrule
\end{tabular}
\end{table}

Table~\ref{tab:meta} compares both models with a reference value and with three metaheuristics that search over visiting orders and score every order with the same exact SOCP. The reference value of each instance is the best upper bound over all methods in \citet{li2025}. It is not a best known solution in general, which is why a few entries in the table are negative. \placeholder{Merge the best values of Vieira et al. (2026) into the reference; per-instance values are needed, since their preprint reports the Poikonen set as means by size.} Each metaheuristic was given one, three or ten times the single-core time that E133 needs on that instance. At equal time, iterated local search (ILS) and variable neighbourhood search (VNS) are worse than E133 overall (+2.93\%), but they are already better than both learned policies on the clustered instances of \citet{poikonen2020} (+1.30\% and +1.31\% against +2.10\% for E133 and +1.76\% for G200). At three times the time their mean gaps (+1.63\% and +1.65\%) are still larger than that of E133, although ILS is better than E133 on more instances than it is worse (251 against 180). At ten times the time they are better than both learned policies, at +0.73\% and +0.78\% against +1.35\% for E133 and +1.14\% for G200. Clustered instances differ from the uniform distribution used for training, which is consistent with the weakness of both learned policies there. Simulated annealing (SA) remains behind at all budgets, and its gap hardly changes between one and three times the time (+4.53\% and +4.52\%), which indicates that our SA implementation is poorly calibrated for this problem; we report its ten-times row for completeness and draw no conclusion from it. On average, E133 and G200 need 0.78\,s of wall time per instance, or 10.3\,s on a single core, on the same hardware. The time grows with size, from under 0.1\,s of wall time at $J = 10$ to about 5.6\,s at $J = 200$, with 16 parallel solver workers.

\zh{表~\ref{tab:meta} 将两个模型与参考值以及三种元启发式进行比较，这些元启发式在访问顺序上搜索，并用同一个精确 SOCP 为每个顺序打分。每个实例的参考值是 \citet{li2025} 中所有方法的最佳上界。它一般不是已知最优解，因此表中少数条目为负。\placeholder{将 Vieira 等（2026）的最佳值并入参考值；需要逐实例的数值，因为其预印本对 Poikonen 集只按规模报告均值。}每种元启发式被给予 E133 在该实例上所需单核时间的一倍、三倍或十倍。在同等时间下，迭代局部搜索（ILS）和变邻域搜索（VNS）总体上不如 E133（+2.93\%），但在 \citet{poikonen2020} 的聚类实例上已经优于两个学习型策略（+1.30\% 和 +1.31\%，对比 E133 的 +2.10\% 和 G200 的 +1.76\%）。在三倍时间下，它们的平均差距（+1.63\% 和 +1.65\%）仍大于 E133，不过 ILS 优于 E133 的实例多于劣于 E133 的实例（251 对 180）。在十倍时间下，它们优于两个学习型策略，分别为 +0.73\% 和 +0.78\%，而 E133 为 +1.35\%、G200 为 +1.14\%。聚类实例不同于训练所用的均匀分布，这与两个学习型策略在其上的弱势一致。模拟退火（SA）在所有预算下都落后，且在一倍和三倍时间之间差距几乎不变（+4.53\% 和 +4.52\%），这表明我们的 SA 实现对该问题校准不佳；我们出于完整性报告其十倍时间的结果，但不据此得出结论。平均而言，在同一硬件上，E133 和 G200 每个实例需要 0.78\,s 墙钟时间，或 10.3\,s 单核时间。时间随规模增长，在 16 个并行求解进程下，从 $J = 10$ 时不到 0.1\,s 增加到 $J = 200$ 时约 5.6\,s 墙钟时间。}

The subset averages in Table~\ref{tab:meta} hide a strong effect of size. Measured against the same reference, E133 and G200 are +4.50\% and +4.02\% above it at $J = 100$, and +5.91\% and +7.09\% at $J = 200$. The two-stage iterative heuristic (TSIH) of \citet{li2025} is +0.30\% above the reference over all 529 instances at a mean of 20.2\,s per instance, and +0.51\% and +0.46\% at $J = 100$ and $J = 200$, at 45.9\,s and 164\,s. These figures are taken from \citet{li2025}, whose hardware differs from ours, and the times are not normalised. The learned policies are therefore far from the strongest published heuristic at the largest sizes, which is also where their advantage in wall time is largest. \placeholder{After the Vieira merge, count instances on which a method in this paper is below the reference by more than rounding (reference values have two decimals), recheck each with an independent feasibility check, and report them. No new best solutions are claimed until then.}

\zh{表~\ref{tab:meta} 中的子集平均掩盖了很强的规模效应。以同一参考值衡量，E133 和 G200 在 $J = 100$ 时分别比参考值高 +4.50\% 和 +4.02\%，在 $J = 200$ 时分别高 +5.91\% 和 +7.09\%。\citet{li2025} 的两阶段迭代启发式（TSIH）在全部 529 个实例上比参考值高 +0.30\%，平均每个实例 20.2\,s；在 $J = 100$ 和 $J = 200$ 时为 +0.51\% 和 +0.46\%，用时 45.9\,s 和 164\,s。这些数据取自 \citet{li2025}，其硬件与我们不同，时间也未归一化。因此，学习型策略在最大规模上与最强的已发表启发式相距甚远，而这恰恰也是它们墙钟时间优势最大的地方。\placeholder{合并 Vieira 的结果后，统计本文方法比参考值低出超过舍入误差的实例（参考值保留两位小数），对每个实例做独立的可行性复核并报告。在此之前不声称任何新的最优解。}}

\begin{table}[!htbp]
\centering
\caption{Real529 gap to the reference value $R$ (best upper bound over all methods in \citet{li2025}), $100\,(\mathrm{UB}-R)/\mathrm{UB}$, where UB is the upper bound found by the method, by subset. Budgets are multiples of the single-core time of E133 on each instance. Our SA implementation is poorly calibrated (see text).\protect\newline\zhcap{Real529 上相对参考值 $R$（\citet{li2025} 中所有方法的最佳上界）的差距 $100\,(\mathrm{UB}-R)/\mathrm{UB}$，其中 UB 为方法找到的上界，按子集列出。预算为 E133 在每个实例上单核时间的倍数。我们的 SA 实现校准不佳（见正文）。}}
\label{tab:meta}
\small
\resizebox{\textwidth}{!}{\begin{tabular}{lrrrrrr}
\toprule
Method & Gam72 & \begin{tabular}[c]{@{}c@{}}Erd132\\ 16--20\end{tabular} & \begin{tabular}[c]{@{}c@{}}Erd132\\ 25--50\end{tabular} & \begin{tabular}[c]{@{}c@{}}Poi325\\ Regular\end{tabular} & \begin{tabular}[c]{@{}c@{}}Poi325\\ Cluster\end{tabular} & All \\
\midrule
E133 & +0.19 & +0.24 & +0.89 & +1.75 & +2.10 & +1.35 \\
G200 & +0.01 & +0.08 & +0.80 & +1.59 & +1.76 & +1.14 \\
ILS, $1\times$ & +1.20 & +1.40 & +4.21 & +5.03 & +1.30 & +2.93 \\
VNS, $1\times$ & +1.20 & +1.40 & +4.21 & +5.04 & +1.31 & +2.93 \\
ILS, $3\times$ & +0.71 & +0.80 & +3.08 & +2.79 & +0.36 & +1.63 \\
VNS, $3\times$ & +0.69 & +0.76 & +3.14 & +2.84 & +0.36 & +1.65 \\
ILS, $10\times$ & +0.33 & +0.27 & +1.90 & +1.22 & $-$0.02 & +0.73 \\
VNS, $10\times$ & +0.20 & +0.34 & +1.89 & +1.40 & $-$0.03 & +0.78 \\
SA, $10\times$ & +0.53 & +1.27 & +5.15 & +5.16 & +1.84 & +3.14 \\
\bottomrule
\end{tabular}}
\end{table}

\subsection{Results on synthetic instances\quad 合成实例上的结果}
\label{sec:synthetic}

Table~\ref{tab:test500} reports test500 with eight views. E133 is behind G200 at every size, with a mean per-instance gap of +1.13\% [+1.02, +1.25]. The gap is smallest at $J = 20$ (+0.46\%), rises to +1.48\% at $J = 80$ and is +1.39\% at $J = 100$. E133 wins 66 instances, ties 51 and loses 383. Most ties occur at $J = 20$, whereas G200 wins at least 82 of 100 instances at each larger size. The shortfall is therefore consistent across the training size range, rather than confined to a few difficult instances. The mean costs are 17.400 for E133 and 17.219 for G200, giving a ratio-of-means gap of +1.05\%, distinct from the mean per-instance gap.

\zh{表~\ref{tab:test500} 报告 test500 的八视角结果。E133 在每个规模上都落后于 G200，逐实例平均差距为 +1.13\% [+1.02, +1.25]。差距在 $J = 20$ 时最小（+0.46\%），在 $J = 80$ 时增至 +1.48\%，在 $J = 100$ 时为 +1.39\%。E133 胜 66、平 51、负 383。大多数平局出现在 $J = 20$，而在每个更大规模上，G200 每 100 个实例至少赢 82 个。因此，这种落后贯穿训练规模范围，而不是集中在少量困难实例上。E133 和 G200 的平均代价分别为 17.400 和 17.219，均值之比对应的差距为 +1.05\%，不同于逐实例平均差距。}

\begin{table}[!htbp]
\centering
\caption{Test500 with eight-view augmentation. Gap is E133 relative to G200 (positive means E133 is worse).\protect\newline\zhcap{Test500 的八视角增强结果。差距为 E133 相对 G200（正值表示 E133 更差）。}}
\label{tab:test500}
\small
\begin{tabular}{rrrrrr}
\toprule
$J$ & E133 & G200 & Gap (\%) & 95\% interval & Win/Tie/Loss \\
\midrule
20 & 12.464 & 12.416 & +0.46 & [+0.29, +0.63] & 8/45/47 \\
40 & 15.634 & 15.500 & +1.02 & [+0.78, +1.26] & 14/4/82 \\
60 & 17.988 & 17.782 & +1.32 & [+1.06, +1.57] & 14/0/86 \\
80 & 19.491 & 19.242 & +1.48 & [+1.19, +1.77] & 17/0/83 \\
100 & 21.422 & 21.154 & +1.39 & [+1.10, +1.67] & 13/2/85 \\
All & 17.400 & 17.219 & +1.13 & [+1.02, +1.25] & 66/51/383 \\
\bottomrule
\end{tabular}
\end{table}

\subsection{Generalisation to larger instances\quad 向更大规模实例的泛化}

Table~\ref{tab:gen} reports instances with 150 to 500 targets, which are beyond the training range, under the same eight-view protocol. The mean per-instance gap to G200 is +0.72\% [+0.53, +0.90], and at $J = 300$ the interval includes zero. The mean costs are 35.827 for E133 and 35.644 for G200; E133 wins 135 instances and loses 265. Neither model was trained at these sizes. Because the two models differ in network and training length as well as in training signal, these results describe the two recipes; they do not show what the training signal alone contributes to transfer.

\zh{表~\ref{tab:gen} 报告训练范围之外的 150 至 500 个目标实例，采用相同的八视角协议。与 G200 的逐实例平均差距为 +0.72\% [+0.53, +0.90]，在 $J = 300$ 时区间包含零。E133 和 G200 的平均代价分别为 35.827 和 35.644；E133 胜 135、负 265。两个模型都没有在这些规模上训练过。由于两个模型除训练信号外在网络和训练长度上也不同，这些结果描述的是两种训练方案，并不说明训练信号本身对泛化的贡献。}

\begin{table}[!htbp]
\centering
\caption{Generalisation set with eight views. Gap is E133 relative to G200.\protect\newline\zhcap{八视角下的泛化集。差距为 E133 相对 G200。}}
\label{tab:gen}
\small
\begin{tabular}{rrrrrr}
\toprule
$J$ & E133 & G200 & Gap (\%) & 95\% interval & Win/Tie/Loss \\
\midrule
150 & 26.003 & 25.795 & +0.96 & [+0.62, +1.31] & 23/0/77 \\
200 & 29.077 & 28.868 & +0.97 & [+0.56, +1.37] & 30/0/70 \\
300 & 36.945 & 36.875 & +0.38 & [$-$0.03, +0.78] & 42/0/58 \\
500 & 51.283 & 51.040 & +0.56 & [+0.23, +0.89] & 40/0/60 \\
All & 35.827 & 35.644 & +0.72 & [+0.53, +0.90] & 135/0/265 \\
\bottomrule
\end{tabular}
\end{table}

\subsection{Comparison with a general-purpose exact solver\quad 与通用精确求解器的比较}

To set the scale of the problem, we also solved the full mixed-integer SOCP (Model+ of \citet{li2025}) with Gurobi 12.0.3, with a time limit of 3,600\,s per instance, one thread and a relative optimality tolerance of $10^{-6}$. No learned route was supplied as a warm start. This is a comparison with a general-purpose solver on the full visiting-order problem, not with the fixed-order SOCP used to evaluate policy routes.

\zh{为说明问题的求解难度，我们还用 Gurobi 12.0.3 求解完整的混合整数 SOCP（\citet{li2025} 的 Model+），每个实例限时 3,600\,s，使用单线程，相对最优性容差为 $10^{-6}$。没有将学习策略的路线作为热启动。这比较的是通用求解器对完整访问顺序问题的求解，而不是用于评估策略路线的定序 SOCP。}

Table~\ref{tab:fullgurobi} reports the results. Gurobi proved 136 of 529 Real529 instances optimal; the other 393 reached the time limit, including all 75 instances with $J \ge 100$, for which no incumbent was found. On test500, six instances were proved optimal, all at $J = 20$, and 494 reached the limit; 49 at $J = 100$ had no incumbent. Upper bounds and relative gaps are averaged over instances with an incumbent, while lower bounds and times cover all instances. Missing incumbents are not assigned an infinite cost.

\zh{表~\ref{tab:fullgurobi} 报告了结果。Gurobi 证明了 Real529 的 529 个实例中 136 个的最优性，其余 393 个达到时间上限，其中包括全部 75 个 $J \ge 100$ 的实例，这些实例均未找到可行解。在 test500 上，六个实例被证明最优，均为 $J = 20$，494 个达到时间上限；$J = 100$ 时有 49 个实例没有找到可行解。上界和相对差距在找到可行解的实例上取均值，下界和时间则覆盖全部实例。没有找到可行解的实例不被赋予无限代价。}

\begin{table}[!htbp]
\centering
\caption{Full mixed-integer SOCP with a 3,600\,s limit and one thread. Opt. is the number proved optimal; Inc. is the number with a feasible incumbent. UB and relative solver gap $100(\mathrm{UB}-\mathrm{LB})/|\mathrm{UB}|$ are means over Inc. only. LB and solve time are means over all $n$ instances. A dash means no incumbent was found.\protect\newline\zhcap{完整混合整数 SOCP，限时 3,600\,s，使用单线程。Opt. 为证明最优的实例数，Inc. 为找到可行解的实例数。UB 和相对求解器差距 $100(\mathrm{UB}-\mathrm{LB})/|\mathrm{UB}|$ 仅在 Inc. 个实例上取均值，LB 和求解时间在全部 $n$ 个实例上取均值。横线表示没有找到可行解。}}
\label{tab:fullgurobi}
\small
\resizebox{\textwidth}{!}{\begin{tabular}{lrrrrrrr}
\toprule
Set and size & $n$ & Opt. & Inc. & UB & LB & Gap (\%) & Time (s) \\
\midrule
Real529, $J = 10$--19 & 220 & 132 & 220 & 254.824 & 236.612 & 7.36 & 1,829.8 \\
Real529, $J = 20$--50 & 234 & 4 & 234 & 370.943 & 260.101 & 29.41 & 3,568.8 \\
Real529, $J = 100$ & 50 & 0 & 0 & -- & 303.108 & -- & 3,600.1 \\
Real529, $J = 200$ & 25 & 0 & 0 & -- & 511.578 & -- & 3,600.0 \\
Real529, all & 529 & 136 & 454 & 314.674 & 266.282 & 18.73 & 2,850.0 \\
\midrule
Test500, $J = 20$ & 100 & 6 & 100 & 12.430 & 8.108 & 37.61 & 3,418.1 \\
Test500, $J = 40$ & 100 & 0 & 100 & 17.099 & 9.143 & 47.56 & 3,600.0 \\
Test500, $J = 60$ & 100 & 0 & 100 & 23.181 & 10.497 & 55.19 & 3,600.1 \\
Test500, $J = 80$ & 100 & 0 & 100 & 31.189 & 11.250 & 63.30 & 3,601.1 \\
Test500, $J = 100$ & 100 & 0 & 51 & 650.174 & 12.110 & 70.98 & 3,607.3 \\
Test500, all & 500 & 6 & 451 & 92.126 & 10.222 & 53.18 & 3,565.3 \\
\bottomrule
\end{tabular}}
\end{table}

On the matched subsets with an incumbent, the mean per-instance gap of E133 to the Gurobi incumbent is $-3.18$\% on 454 Real529 instances and $-20.25$\% on 451 test500 instances. The large test500 difference reflects weak incumbents at the time limit, especially at $J = 100$, where their mean cost is 650.174 on the 51 instances with an incumbent; it is not evidence that E133 is close to the optimum. On the subsets proved optimal, E133 is 0.26\% above the optimum on 136 Real529 instances and 0.10\% above it on six test500 instances. Those subsets contain small instances and do not establish optimality gaps for the larger ones. The learned policy supplies a feasible solution for every instance, but neither its advantage over a time-limited incumbent nor its exact fixed-order evaluation is a certificate of global optimality.

\zh{在找到可行解的配对子集上，E133 相对 Gurobi 当前最好解的逐实例平均差距，在 Real529 的 454 个实例上为 $-3.18$\%，在 test500 的 451 个实例上为 $-20.25$\%。test500 上的较大差异反映了时间上限时解的质量较差，尤其是 $J = 100$ 时，51 个找到可行解的实例上平均代价为 650.174；这并不表明 E133 接近最优解。在证明最优的子集上，E133 在 Real529 的 136 个实例上比最优值高 0.26\%，在 test500 的六个实例上高 0.10\%。这些子集包含的是小规模实例，不能据此确定更大实例的最优性差距。学习策略为每个实例都给出了可行解，但它相对于限时可行解的优势，以及其定序评估的精确性，都不是全局最优性的证书。}

\subsection{Training efficiency\quad 训练效率}
\label{sec:efficiency}

We compare training budgets in two ways. The first holds the solver budget fixed. POMO trained from scratch under the G200 recipe on our workstation was stopped once it had used the same 3.84 million exact solves as E133, which took it 264,512 instances. With eight-view augmentation, its mean per-instance gap to E133 is +3.79\% [+3.28, +4.30] on Real529 and +3.27\% [+2.89, +3.65] on test500. E133 is better on 353 Real529 instances, tied on 111 and worse on 65; on test500, the counts are 369, 37 and 94. The Real529 gap grows with size, to 11.7\% at $J = 100$ and 24.0\% at $J = 200$. These results compare two recipes under a small solver budget, not the final quality of exact-reward POMO. The second comparison is with G200 at its full budget, which remains better than E133 overall on both sets (Sections~\ref{sec:real} and~\ref{sec:synthetic}).

\zh{我们以两种方式比较训练预算。第一种固定求解器预算。在我们的工作站上按 G200 方案从头训练的 POMO，在用完与 E133 相同的 384 万次精确求解时停止，此时它用了 264,512 个实例。采用八视角增强时，它相对 E133 的逐实例平均差距在 Real529 上为 +3.79\% [+3.28, +4.30]，在 test500 上为 +3.27\% [+2.89, +3.65]。E133 在 Real529 上胜 353、平 111、负 65；在 test500 上相应为 369、37 和 94。Real529 上的差距随规模增大，在 $J = 100$ 和 $J = 200$ 时分别为 11.7\% 和 24.0\%。这些结果比较的是小求解器预算下的两套方案，而不是精确奖励 POMO 的最终质量。第二种比较是与完整预算下的 G200 比较，它在两个集合的总体结果上仍优于 E133（第~\ref{sec:real} 节和第~\ref{sec:synthetic} 节）。}

Table~\ref{tab:budget} lists the full budgets. E133 used 15 times fewer training instances and 78 times fewer Gurobi calls in training than G200. The second ratio is the product of two factors: 15 times fewer instances and 5.2 times fewer calls per instance. G200 solved $J/4$ routes per training instance, so a baseline that trained with $J$ starts, as both models use at test time, would have needed about four times as many calls. These ratios stop at E133 and count training only. The run that produced E133 continued to epoch 200, using 6,239,185 training calls and 400,000 validation calls, and selecting E133 among its checkpoints by full POMO validation took 5,175,733 further solves. As run, the total is therefore 11,814,918 calls. The training log of G200 counts 338,545,920 calls when the epochs that were restarted after interruptions are counted each time they ran, about 29 times our total. G200 used its final epoch and needed no checkpoint selection, and its log records no validation calls. The test-time cost is unchanged in our measurements because it is dominated by exact solves of the POMO routes: on Real529, decoding takes 0.12\,s per instance on the GPU and Gurobi takes about 10.2\,s of single-core time for either model.

\zh{表~\ref{tab:budget} 列出了完整预算。E133 使用的训练实例比 G200 少 15 倍，训练中的 Gurobi 调用少 78 倍。第二个比值是两个因子的乘积：实例少 15 倍，每个实例的调用少 5.2 倍。G200 对每个训练实例求解 $J/4$ 条路线，因此一个像两个模型在测试时那样以 $J$ 个起点训练的基线，大约需要四倍的调用。这些比值截止于 E133，且只计训练。产生 E133 的运行一直持续到第 200 轮，使用了 6,239,185 次训练调用和 400,000 次验证调用，而用完整 POMO 验证从其检查点中选出 E133 又用了 5,175,733 次求解。因此按实际运行计，总量为 11,814,918 次调用。G200 的训练日志中，若把中断后重跑的轮次按实际运行次数重复计入，共计 338,545,920 次调用，约为我们总量的 29 倍。G200 使用其最后一轮，无需检查点选择，其日志也没有记录验证调用。在我们的测量中，测试时成本不变，因为它主要由 POMO 路线的精确求解决定：在 Real529 上，GPU 解码每个实例需 0.12\,s，Gurobi 对任一模型约需 10.2\,s 单核时间。}

The logged single-thread solver runtimes of E133 sum to 59.9 hours, or 0.056\,s per call over its mix of sizes. Initial table search adds 67.6 hours of summed worker time. A further logged 26.9 hours combines search and solving inside the refresh rounds; its solver component is already included in the 59.9 hours and cannot be separated from the saved timing columns. These entries therefore cannot be added without double counting. The resulting solver-plus-search time lies between 127.5 and 154.4 summed worker-hours, excluding GPU work and unlogged overhead; it is not a measurement of process CPU time. The solver runtime of G200 is only an estimate, obtained by multiplying its call count by a mean solve time of 0.031\,s measured at about $J = 60$. Its wall time also comes from a different machine with 128 reward workers, so the ratio of wall times shows the order of magnitude only. Solver calls and instances are exact counts and form the primary comparison.

\zh{E133 日志中的单线程求解器运行时间合计为 59.9 小时，在其规模组合上平均每次调用为 0.056\,s。初始表搜索另有 67.6 小时的工作进程累计时间。日志还记录了 26.9 小时的刷新轮次搜索与求解混合时间，其中求解部分已包含在 59.9 小时中，且无法从保存的计时列中单独分离。因此，直接相加会重复计数。由此得到的求解器加搜索时间介于 127.5 和 154.4 个累计工作进程小时之间，不含 GPU 工作和未记录的开销；它不是进程 CPU 时间的测量。G200 的求解器运行时间只是估计值，由其调用次数乘以在 $J \approx 60$ 附近测得的平均求解时间 0.031\,s 得到。其墙钟时间也来自另一台配有 128 个奖励进程的机器，因此墙钟时间之比只表示数量级。求解器调用次数和实例数是精确计数，构成主要比较。}

Table~\ref{tab:calls} breaks down the 3.84 million calls of E133. Anchor and target solves contribute one call per instance each. Expert verification, including near-miss checks, contributes 0.37 calls per instance. Refresh rounds add 0.52 calls per instance on average across all sizes, and more than one per instance at the end of the run, when more than half of the batches have $J \ge 60$ and refresh is active. A share of 23.7\% of refresh solves produced a better expert, so the refresh budget has a measurable return but also a large waste. Near-miss checks were attempted 364,974 times.

\zh{表~\ref{tab:calls} 分解了 E133 的 384 万次调用。锚点求解和目标求解各贡献每个实例一次调用。专家验证（包括近失检验）贡献每个实例 0.37 次。刷新轮次在所有规模上平均每个实例增加 0.52 次，在运行末期超过每个实例一次，此时超过一半的批次 $J \ge 60$、刷新处于激活状态。刷新求解中有 23.7\% 产生了更好的专家，因此刷新预算有可衡量的回报，但也有很大浪费。近失检验共尝试了 364,974 次。}

\begin{table}[!htbp]
\centering
\caption{Exact solver calls of E133, epochs 1 to 133, by purpose.\protect\newline\zhcap{E133 第 1 至 133 轮按用途划分的精确求解器调用。}}
\label{tab:calls}
\small
\begin{tabular}{lrr}
\toprule
Purpose & Calls & Per instance \\
\midrule
Anchor solve & 1,330,000 & 1.00 \\
Target reward solve & 1,330,000 & 1.00 \\
Expert and near-miss verification & 493,128 & 0.37 \\
Refresh rounds & 686,842 & 0.52 \\
Total & 3,839,970 & 2.89 \\
\bottomrule
\end{tabular}
\end{table}

Table~\ref{tab:dynamics} shows how the loop behaves during training. In the first epoch, almost any search improves on an untrained policy. By epoch 40 the plain 2-opt search improves only 9\% of the anchors, by 0.38\% on average. The ramp then lifts the share of improved experts above one third, which is the intended effect of the progressive search. Over the same period, the lowest-bound policy route has a bound that falls from 92\% to 34\% above its exact cost, which is consistent with policy routes moving closer to the anchors as training proceeds. The bound remains loose at epoch 133, yet training continues to make progress; Section~\ref{sec:loose} discusses what this does and does not show.

\zh{表~\ref{tab:dynamics} 展示了训练过程中循环的行为。在第一轮，几乎任何搜索都能改进未经训练的策略。到第 40 轮，普通 2-opt 搜索只改进了 9\% 的锚点，平均改进 0.38\%。之后的斜坡把改进专家的比例提高到三分之一以上，这正是渐进式搜索的预期效果。在同一时期，策略路线中最低上界相对于该路线精确代价的超出量从 92\% 降到 34\%，这与训练中策略路线向锚点靠拢一致。上界在第 133 轮仍然宽松，但训练仍在进步；第~\ref{sec:loose} 节讨论这说明了什么、没有说明什么。}

\begin{table}[!htbp]
\centering
\caption{Training dynamics of the proposed method. Improved share is the fraction of instances whose expert beats the anchor. Bound excess compares the lowest bound among policy routes with the exact cost of that route.\protect\newline\zhcap{所提方法的训练动态。改进比例是专家优于锚点的实例比例。上界超出量比较策略路线中的最低上界与该路线的精确代价。}}
\label{tab:dynamics}
\small
\begin{tabular}{rrrrr}
\toprule
Epoch & Anchor cost & Expert cost & Improved share (\%) & Bound excess (\%) \\
\midrule
20 & 19.986 & 19.779 & 12.1 & 92.4 \\
40 & 19.205 & 19.096 & 9.0 & 70.7 \\
60 & 19.146 & 18.798 & 24.4 & 57.5 \\
80 & 18.993 & 18.640 & 33.0 & 57.9 \\
100 & 18.850 & 18.603 & 34.9 & 44.2 \\
133 & 18.522 & 18.288 & 37.9 & 33.6 \\
\bottomrule
\end{tabular}
\end{table}

\begin{table}[!htbp]
\centering
\caption{Training budgets of G200 and E133.\protect\newline\zhcap{G200 与 E133 的训练预算。}}
\label{tab:budget}
\small
\resizebox{\textwidth}{!}{\begin{tabular}{lrrr}
\toprule
 & G200 & E133 & Ratio \\
\midrule
Training instances (to E133) & 20,000,000 & 1,330,000 & 15 \\
Gurobi calls in training (to E133) & 299,437,760 & 3,839,970 & 78 \\
Gurobi calls as run\textsuperscript{a} & 338,545,920 & 11,814,918 & 29 \\
Summed solver runtime (h, to E133) & about 2,580 (estimated at 0.031\,s per call) & 59.9 (0.056\,s per call) & -- \\
Initial table-search worker time (h, to E133) & -- & 67.6 & -- \\
Refresh search-and-solve worker time (h, to E133)\textsuperscript{b} & -- & 26.9 & -- \\
Wall time (hours) & about 215.5 (128 workers, other machine) & 19.2 to E133, 28.3 to E200 (32 workers) & about 11 \\
Test time per Real529 instance (s) & 0.78 & 0.78 & 1 \\
\bottomrule
\end{tabular}}
\par\smallskip\footnotesize\raggedright \textsuperscript{a}G200: all training epochs, with epochs restarted after interruptions counted each time they ran; it used its final epoch, so no checkpoint selection, and its log records no validation calls. E133: training to E200 (6,239,185), in-training validation (400,000) and checkpoint selection on validation2000 (5,175,733).
\par \textsuperscript{b}Refresh solver time is also included in the solver-runtime row. The three runtime rows are not additive; the saved logs do not separate refresh search from refresh solves.
\end{table}

Group A repeats the exact-reward POMO training on the same workstation with a 24-hour wall-time budget. Its saved checkpoint falls during epoch 37, after 53,267,040 training solves and 3,654,784 instances. This is a separate training run from the equal-solve comparison above. No eight-view evaluation of its checkpoints is available in the saved results, so we report its training budget only and make no claim about its solution quality here.

\zh{A 组在同一工作站上重复精确奖励 POMO 的训练，墙钟时间预算为 24 小时。保存的检查点位于第 37 轮进行中，已使用 53,267,040 次训练求解和 3,654,784 个实例。这与前面的等求解次数比较是一次独立的训练运行。保存的结果中没有其检查点的八视角评估，因此此处只报告训练预算，不对其解质量作判断。}

\subsection{Ablation study\quad 消融实验}
\label{sec:ablation}

The ablations measure what individual components contribute at a fixed point in training; they report component effects and are not offered as a decomposition of the saving. Variants B and D train from scratch to one million instances (epoch 100), with the configuration and seed of the proposed method otherwise unchanged. They are compared with Group C, the proposed method at the same point, on the same 500 test instances with eight views, by size. We fixed the decision rule before seeing the results: an effect is reported only if the paired per-instance relative cost gap between the variant and C has a 95\% interval that excludes zero. The intervals are pointwise and are not adjusted for comparisons across sizes. Conclusions apply to epoch 100, before the reported checkpoint E133, and each variant is a single run. The comparison fixes training instances and epochs, not solver calls or wall time.
\zh{消融实验衡量各部件在一个固定训练时点的贡献，报告部件效应，而不把节省量分解归因。变体 B 和 D 从头训练至一百万个实例（第 100 轮），除此之外沿用所提方法的配置和种子。它们与同一时点的 C 组（所提方法）在同一批 500 个测试实例上使用八视角、按规模比较。我们在看到结果前固定了判定规则：只有变体相对于 C 的逐实例配对相对代价差距的 95\% 区间不包含零时，才报告存在效应。这些区间是逐项区间，未对跨规模比较做调整。结论适用于第 100 轮，早于报告的检查点 E133，且每个变体都只有一次运行。比较固定的是训练实例数和轮次，不是求解器调用次数或墙钟时间。}
\begin{itemize}
\item Group B lets the expert search evaluate candidate moves exactly instead of on the table. Because exact evaluation is expensive, the search is also reduced: it starts from the anchor only, runs first-improvement 2-opt with at most 100 exact solves per instance, and uses no LKH and no refresh. The reward for policy routes still uses the anchor table. Group B therefore measures the role of the table inside the expert search, together with the effect of a smaller search, and not its role as a reward. To epoch 100 it made 102,762,944 exact solves, including 100 million exact neighbourhood solves and one million target-reward solves, against 2,643,875 for C including refresh. This is 38.9 times as many calls, and every instance used its full search budget. Its training took 4.4 times as long (48.6 against 11.0 hours), excluding validation time from both runs. \zhitem{B 组让专家搜索精确评估候选移动，而不是用表评估。由于精确评估昂贵，搜索也被缩小：只从锚点出发，运行首次改进的 2-opt，每个实例最多做 100 次精确求解，不使用 LKH 或刷新。策略路线的奖励仍使用锚点表。因此 B 组衡量的是表在专家搜索中的作用以及搜索缩小的共同效应，而不是表作为奖励的作用。到第 100 轮，它使用了 102,762,944 次精确求解，包括一亿次邻域精确求解和一百万次目标奖励求解；C 组含刷新为 2,643,875 次。这是 38.9 倍的调用，每个实例都用满了搜索预算。其训练时间为 C 的 4.4 倍（48.6 对 11.0 小时），两边都不含验证时间。}
\item Group D reduces the reference pool to the greedy route of the policy from the first start, which receives the local-search budget of the 30 starts. It varies the diversity of the starts of the expert search, together with their quality. It does not vary the diversity of the POMO routes discussed in Section~\ref{sec:limits}. \zhitem{D 组把参考池缩减为策略从第一个起点出发的贪心路线，并让它获得 30 个起点的局部搜索预算。它改变的是专家搜索起点的多样性，同时也改变了起点的质量；它并不改变第~\ref{sec:limits} 节讨论的 POMO 路线多样性。}
\end{itemize}

Table~\ref{tab:ablationBCD} reports the completed evaluations. Group B has a mean per-instance relative gap to C of $-0.10$\% [$-0.21$, $+0.02$], so the overall interval includes zero. Its gains at $J = 20$ and $J = 40$ have intervals below zero, but those at the three larger sizes do not. B wins 251 instances, ties 62 and loses 187. Its arithmetic mean cost is slightly higher than C's (17.464 against 17.454), even though its mean relative gap is negative, because the modest gains on smaller instances are offset by larger absolute costs at $J = 80$ and $J = 100$. Under the stated rule, the much larger solve budget does not establish an overall quality improvement. This is not an equivalence result, and the reduced exact search does not isolate the table's effect from the other search changes.

\zh{表~\ref{tab:ablationBCD} 报告了已完成的评估。B 组相对 C 的逐实例平均相对差距为 $-0.10$\% [$-0.21$, $+0.02$]，总体区间包含零。它在 $J = 20$ 和 $J = 40$ 时的改善区间低于零，但在更大三个规模上不是如此。B 在 251 个实例上胜、62 个持平、187 个负。尽管平均相对差距为负，其算术平均代价仍略高于 C（17.464 对 17.454），因为小实例上的少量改善被 $J = 80$ 和 $J = 100$ 时更大的绝对代价抵消。按照既定规则，大幅增加求解预算并未确立总体质量改善。这不是等价性结果，缩小后的精确搜索也不能把表的效应与其他搜索变化分开。}

\begin{table}[!htbp]
\centering
\caption{Ablations at epoch 100 on test500 with eight views, 100 instances per size. Costs are means. Gaps are mean per-instance percentages relative to C, with 95\% intervals; positive values mean that the variant is worse. B uses exact expert search, C is the proposed recipe, and D uses a single greedy search start.\protect\newline\zhcap{第 100 轮在 test500 上的八视角消融，每个规模 100 个实例。代价为均值，差距为相对 C 的逐实例平均百分比及其 95\% 区间，正值表示变体更差。B 使用精确专家搜索，C 为所提方案，D 使用单条贪心搜索起点。}}
\label{tab:ablationBCD}
\small
\resizebox{\textwidth}{!}{\begin{tabular}{rrrrrrrr}
\toprule
$J$ & C cost & B cost & B gap (\%) & 95\% interval & D cost & D gap (\%) & 95\% interval \\
\midrule
20 & 12.480 & 12.449 & $-$0.31 & [$-$0.46, $-$0.16] & 12.501 & +0.20 & [+0.03, +0.38] \\
40 & 15.646 & 15.622 & $-$0.27 & [$-$0.49, $-$0.06] & 15.850 & +1.33 & [+1.00, +1.66] \\
60 & 18.017 & 17.999 & $-$0.21 & [$-$0.47, +0.05] & 18.427 & +2.36 & [+1.95, +2.77] \\
80 & 19.562 & 19.611 & +0.13 & [$-$0.13, +0.39] & 20.178 & +3.21 & [+2.81, +3.61] \\
100 & 21.568 & 21.637 & +0.19 & [$-$0.15, +0.52] & 22.334 & +3.59 & [+3.18, +3.99] \\
All & 17.454 & 17.464 & $-$0.10 & [$-$0.21, +0.02] & 17.858 & +2.14 & [+1.94, +2.33] \\
\bottomrule
\end{tabular}}
\end{table}

Group D is worse than C at every size, with an overall gap of +2.14\% [+1.94, +2.33] and 61 wins, 44 ties and 395 losses. The gap rises from +0.20\% at $J = 20$ to +3.59\% at $J = 100$. It uses fewer training solves, 2,111,419 including 99,069 expert verifications and 12,350 refresh solves; each of B, C and D also incurs 200,000 in-training validation solves, counted separately. The result supports retaining the sampled multistart expert search in this recipe, especially on large instances, but it does not distinguish start quality from diversity. Nor does it test the diversity of the learning policy's POMO output. Across both ablations, the evidence concerns these training recipes at one million instances and one seed, not isolated causal effects or variability across training runs.

\zh{D 组在每个规模上都比 C 差，总体差距为 +2.14\% [+1.94, +2.33]，胜 61、平 44、负 395。差距从 $J = 20$ 时的 +0.20\% 增至 $J = 100$ 时的 +3.59\%。其训练求解次数更少，为 2,111,419 次，其中包括 99,069 次专家验证和 12,350 次刷新求解；B、C、D 各另有 200,000 次训练期间的验证求解，单独计数。结果支持在本方案中保留采样式多起点专家搜索，尤其是对大实例，但不能区分起点质量与多样性，也没有测试学习策略 POMO 输出的多样性。两项消融的证据都针对一百万实例、单个种子下的这些训练方案，而非孤立的因果效应或训练运行之间的变异。}

\section{Discussion and implications for practice\quad 讨论与实践启示}
\label{sec:discussion}

\subsection{When the approach is worthwhile\quad 何时值得采用该方法}

A trained policy moves most of the computing effort into training, but not all of it: at test time E133 still solves one SOCP per distinct decoded route, about 10.3\,s of single-core time per Real529 instance, or 0.78\,s of wall time with 16 solver workers. Within the range of speeds and endurance seen in training, the policy can be reused without retraining, because both enter as input features. The same holds for G200, which gives better plans at the same test time. The smaller training budget therefore matters when a learned solver has to be trained from scratch, for example for a variant with additional constraints or a different spatial distribution of targets. In that situation the solver budget of training decides whether a learned solver is affordable at all, and our results show that a useful policy can be obtained with a few million solves instead of hundreds of millions. We did not test training on a changed problem.

\zh{训练好的策略把大部分计算工作移到训练阶段，但并非全部：在测试时，E133 仍要为每条不同的解码路线求解一个 SOCP，每个 Real529 实例约需 10.3\,s 单核时间，或在 16 个求解进程下 0.78\,s 墙钟时间。在训练中见过的速度和续航范围内，策略无需重新训练即可复用，因为两者都作为输入特征。G200 同样如此，并且在相同测试时间下给出更好的方案。因此，更小的训练预算只在必须从头训练学习型求解器时才有意义，例如面对增加了约束或目标空间分布不同的变体。在这种情况下，训练的求解器预算决定了学习型求解器是否负担得起；我们的结果表明，用几百万次而非几亿次求解即可得到一个有用的策略。我们没有测试在改变后的问题上训练。}

Table~\ref{tab:breakeven} sets the computing cost of the options against their plan quality on Real529. The training CPU of E133 to epoch 133, 154.3 core-hours of Gurobi and table search, equals the additional single-core time of ILS at ten times the budget on about 6,000 instances, at about 93\,s per instance. Below that number of plans, ILS with ten times the budget uses less CPU in total and gives better plans. Above it, the policy uses less CPU, but its plans are about 0.6 percentage points further from the reference on average, and much further at $J \ge 100$ (Section~\ref{sec:real}). Where G200 is already available it is preferable to E133, since its test cost is the same and its plans are better. The comparison treats training CPU and planning CPU as equal in value and ignores GPU time.

\zh{表~\ref{tab:breakeven} 将各方案的计算成本与其在 Real529 上的方案质量对比。E133 训练到第 133 轮的 CPU 时间为 154.3 核时（Gurobi 与表搜索），相当于十倍预算的 ILS 在约 6,000 个实例上的额外单核时间（每个实例约 93\,s）。方案数少于这一数量时，十倍预算的 ILS 总 CPU 更少且方案更好。多于这一数量时，策略使用的 CPU 更少，但其方案平均比参考值远约 0.6 个百分点，在 $J \ge 100$ 时差得更多（第~\ref{sec:real} 节）。若已有 G200，它优于 E133，因为两者测试成本相同而 G200 的方案更好。这一比较把训练 CPU 与规划 CPU 视为同等价值，并忽略了 GPU 时间。}

\begin{table}[!htbp]
\centering
\caption{Computing cost and plan quality on Real529. Single-core time per instance for the metaheuristics is the stated multiple of the mean single-core time of E133. Gap is to the reference value of Table~\ref{tab:meta}.\protect\newline\zhcap{Real529 上的计算成本与方案质量。元启发式每个实例的单核时间为 E133 平均单核时间的相应倍数。差距相对于表~\ref{tab:meta} 的参考值。}}
\label{tab:breakeven}
\small
\resizebox{\textwidth}{!}{\begin{tabular}{lrrrr}
\toprule
Option & Training solver calls & Training CPU (core-hours) & Single-core time per instance (s) & Gap (\%) \\
\midrule
G200 & 299,437,760 & about 2,580 (estimated) & about 10.3 & +1.14 \\
E133 & 3,839,970 (about 11.8 million as run) & 154.3 & about 10.3 & +1.35 \\
ILS, $3\times$ & -- & -- & about 31 & +1.63 \\
ILS, $10\times$ & -- & -- & about 103 & +0.73 \\
\bottomrule
\end{tabular}}
\end{table}

At the time budgets we tested, the advantage of the policy over a metaheuristic is limited. ILS with the same single-core time is behind E133 overall but ahead on clustered instances, and with three times the time it is better on more instances than it is worse (251 against 180), although its mean gap is larger. The clear advantage of the policy lies in wall time when solver workers can run in parallel, which we did not give the metaheuristics. A route from the policy could also start a metaheuristic, and the table could be used for search at test time; we did not test either.

\zh{在我们测试的时间预算下，策略相对于元启发式的优势有限。相同单核时间的 ILS 总体上落后于 E133，但在聚类实例上领先；在三倍时间下，它更好的实例多于更差的实例（251 对 180），尽管其平均差距更大。策略的明显优势在于求解进程可以并行时的墙钟时间，而我们没有让元启发式并行。策略给出的路线还可以作为元启发式的初始解，表也可以用于测试时的搜索；这两点我们都没有测试。}

\subsection{What the loose bound does and does not show\quad 宽松上界说明了什么、没有说明什么}
\label{sec:loose}

The bound is far from tight on routes that differ from the anchor. The lowest bound among policy routes exceeds the exact cost of that route by 33.6\% at epoch 133 (Table~\ref{tab:dynamics}), and near-miss routes, which come out of local search, carry bounds 14.6\% above the anchor cost on average although 40.8\% of them are exactly better (Section~\ref{sec:search}). The bound is thus not accurate near the anchor at the resolution that matters for acceptance, and most late-stage experts are accepted through near-miss checks and refresh rounds. In the recipe it serves mainly as a reward for routes that are never solved and as a ranking that decides which few routes receive an exact solve, while the exactness of every label keeps errors in the bound out of imitation learning.

\zh{对于与锚点不同的路线，上界远不紧。第 133 轮时，策略路线中的最低上界比该路线的精确代价高 33.6\%（表~\ref{tab:dynamics}），而来自局部搜索的近失路线，其上界平均比锚点代价高 14.6\%，尽管其中 40.8\% 精确地更好（第~\ref{sec:search} 节）。因此，在对接受与否起决定作用的精度上，上界在锚点附近并不准确，后期大多数专家是通过近失检验和刷新轮次被接受的。在训练方案中，它主要作为从未求解路线的奖励，以及决定哪几条路线获得精确求解的排序，而每个标签的精确性使上界的误差不会进入模仿学习。}

Two experiments bear on whether a tighter or less biased signal would help, and neither settles the question. In a 30-epoch continuation from E133, taking the minimum over all tables available for each route used about two additional tables per instance on the affected batches and tightened the bound by only 0.32\%. This is a small change against an excess of about one third, so it does not show that tightness has low value. In an earlier 20-epoch study from scratch, rewards from exact solves of a random subset of 4 to 8 routes per instance point in a direction much closer to the all-exact advantage (an expected cosine of about 0.85 to 0.95, against about 0.6 for the hybrid reward). They cost three to four times the Gurobi CPU. These are training-signal diagnostics from short single-seed runs, not eight-view policy-quality comparisons. The hybrid reward is also biased, not only noisy (Section~\ref{sec:learning}), and whether a tighter or unbiased signal would close the remaining gap is open. Two refinements that keep the guarantee are available and untested: letting each target take its sortie from a different solved solution, chosen by a dynamic programme along the route, and letting each target exchange its take-off and landing points, chosen by a two-state dynamic programme along the route.

\zh{有两项实验与更紧或偏差更小的信号是否有帮助相关，但都不能定论。在从 E133 出发的 30 轮续训中，对每条路线取所有可用表上的最小值，在受影响的批次上每个实例约多用两张表，只把上界收紧了 0.32\%。相对于约三分之一的超出量，这是很小的变化，因此不能说明紧度的价值低。在一项更早的、从头开始的 20 轮研究中，对每个实例随机 4 至 8 条路线做精确求解得到的奖励，其方向与全精确优势接近得多（期望余弦约 0.85 至 0.95，而混合奖励约为 0.6）。它们的 Gurobi CPU 成本为三到四倍。这些是短期单种子运行的训练信号诊断，不是八视角策略解质量比较。混合奖励还有偏差，而不仅是有噪声（第~\ref{sec:learning} 节），更紧或无偏的信号能否消除剩余差距仍是开放问题。有两种保持保证的改进可用但尚未测试：让每个目标从不同的已求解方案中取架次（沿路线用动态规划选择），以及让每个目标互换其起飞点和降落点（沿路线用两状态动态规划选择）。}

\subsection{Limitations\quad 局限性}
\label{sec:limits}

Six limitations apply. First, E133 remains behind G200 on synthetic instances from the training distribution, by +1.13\% with eight views. We have not identified the cause, and several explanations remain open. (i) Route diversity: in a diagnostic on a later checkpoint of the same run (epoch 163), the edge overlap between POMO routes was higher than for G200 at every size tested (0.64, 0.67 and 0.62 against 0.48, 0.55 and 0.53 at $J = 20$, 60 and 100). Against this explanation, the overlap difference shrinks with size while the gap grows. (ii) Pull towards the anchor: the bound is tightest at the anchor and the hybrid reward favours the solved route (Section~\ref{sec:learning}), so the reward may draw the policy towards anchor-like routes. In matched 20-epoch runs from the same initialisation, replacing the bound-based reward with random-subset exact rewards lowered the overlap at $J = 20$ from 0.374 to 0.330 and from 0.379 to 0.307. (iii) Training length and recipe: E133 saw 15 times fewer instances and uses a different network, and G200 trained with $J/4$ starts. Continuations from E133 that changed the bound or the imitation weight did not move the overlap appreciably, so they do not discriminate between (i) and (ii). Separating these explanations requires controlled runs that we have not performed. Second, every model is one training run with one seed, and the reported intervals cover variation over test instances only, not over training runs. Large differences, such as the 3.8\% gap at equal solver budget on Real529, are unlikely to reverse with another seed, but small ones, such as the +0.21\% gap on Real529, should be read as closeness and not as a ranking. Third, Real529 is not held out and covers only two parameter settings, both at the edges of the training range. Fourth, the comparison on time depends on hardware and implementation, which is why we treat solver calls as the primary measure. Fifth, the settings of the progressive search were set by development runs and not by a formal search, so the reported performance may not be the best attainable for the method. Sixth, the bound depends on the separable structure of the CVTSP. It extends to a variant only if sortie constraints remain independent of the order, which should hold for several sorties of fixed composition but fails if, for instance, target time windows couple the sorties.

\zh{有六项局限。第一，E133 在训练分布的合成实例上仍落后于 G200，八视角下为 +1.13\%。我们尚未确定原因，有几种解释仍然开放。(i) 路线多样性：在同一运行较晚检查点（第 163 轮）的诊断中，POMO 路线之间的边重合度在所有测试规模上都高于 G200（$J = 20$、60、100 时为 0.64、0.67、0.62，对比 0.48、0.55、0.53）。反对这一解释的证据是：重合度差异随规模缩小，而差距随规模增大。(ii) 向锚点的牵引：上界在锚点处最紧，混合奖励偏向被求解的路线（第~\ref{sec:learning} 节），因此奖励可能把策略拉向与锚点相似的路线。在从相同初始化出发、相互匹配的 20 轮运行中，把基于上界的奖励换成随机子集精确奖励，使 $J = 20$ 时的重合度从 0.374 降到 0.330、从 0.379 降到 0.307。(iii) 训练长度与方案：E133 见过的实例少 15 倍、使用不同的网络，而 G200 以 $J/4$ 个起点训练。从 E133 出发改变上界或模仿权重的续训，没有明显改变重合度，因此它们无法区分 (i) 和 (ii)。要区分这些解释，需要我们尚未进行的受控运行。第二，每个模型都只是一次单种子的训练运行，报告的区间只覆盖测试实例上的变异，不覆盖训练运行之间的变异。较大的差异（如 Real529 上等求解预算下 3.8\% 的差距）不太可能因换种子而逆转，但较小的差异（如 Real529 上 +0.21\% 的差距）应理解为接近，而不是排名。第三，Real529 不是留出集，且只覆盖两种参数设置，都位于训练范围的边缘。第四，时间比较依赖硬件和实现，这就是我们把求解器调用作为主要度量的原因。第五，渐进式搜索的设置来自开发运行，而不是正式搜索，因此报告的性能未必是该方法能达到的最好水平。第六，上界依赖 CVTSP 的可分结构。只有当架次约束仍与顺序无关时，它才能推广到一个变体；对于组成固定的多个架次应当成立，但如果例如目标时间窗把各架次耦合起来，则不成立。}

\section{Conclusion\quad 结论}
\label{sec:conclusion}

This paper asked whether a DRL policy for the CVTSP needs an exact solver reward for every route. One optimal fixed-order solution yields an upper bound on all visiting orders, computed by summing table entries, and a training recipe built on this bound, with 2.89 exact solves per training instance and exact verification of every imitation label, obtains a useful policy at a small fraction of the usual solver budget. With the same 3.84 million solves, POMO trained on exact rewards is on average 3.8\% worse per instance than the policy from our recipe on the literature benchmark. The fully trained baseline used 78 times more solves in training, or about 29 times more when both runs are counted as run, and it remains better: by +0.21\% on the literature benchmark and +0.72\% on larger instances, where the two are close, and by +1.13\% with augmentation on synthetic instances from the training distribution. Iterated local search and variable neighbourhood search given ten times the single-core time of the policy, and the strongest published heuristic, obtain better solutions, most clearly at 100 and 200 targets. The recipe is therefore a way to obtain a learned CVTSP solver with a small solver budget, not a better solver.

\zh{本文研究的问题是：CVTSP 的 DRL 策略是否需要为每条路线提供精确求解器奖励。一个最优定序解给出所有访问顺序的上界，可通过对表元素求和得到；基于这一上界构建的训练方案，每个训练实例使用 2.89 次精确求解，并对每个模仿标签做精确验证，只用通常求解器预算的一小部分就得到了一个有用的策略。在同样 384 万次求解下，以精确奖励训练的 POMO 在文献基准上逐实例平均比本方案的策略差 3.8\%。完全训练的基线在训练中多用了 78 倍的求解，若两边都按实际运行计则约为 29 倍，并且仍然更好：本方案策略与它的差距在文献基准上为 +0.21\%，在更大实例上为 +0.72\%，二者在这些集合上接近；在训练分布的合成实例上，增强后为 +1.13\%。给予策略十倍单核时间的迭代局部搜索和变邻域搜索，以及最强的已发表启发式，都能得到更好的解，在 100 和 200 个目标时最为明显。因此，本方案是以较小求解器预算获得学习型 CVTSP 求解器的方法，而不是一个更好的求解器。}

Several directions follow. The recipe should be tested where training from scratch is needed, such as variants with further constraints or other target distributions. The training signal can be tightened while keeping its guarantee, through the per-target combination of tables and the reversal symmetry. The explanations for the remaining gap, route diversity, the pull of a biased reward towards the anchor and the difference in training length, need to be separated by controlled runs. The table could also support search at test time or warm starts for metaheuristics. The bound requires only that the sorties of one solution remain feasible for other orders, but whether it is useful elsewhere also depends on how loose it is, which we could not predict in advance. Several carriers \citep{fahradyan2020} and several UAVs on one vessel \citep{ren2026} raise different questions, since sorties may then overlap in time and the leg structure changes.

\zh{由此引出几个方向。应在需要从头训练的场景中测试该方案，例如有更多约束或其他目标分布的变体。可以在保持保证的前提下收紧训练信号，方法包括多张表的逐目标组合以及反转对称性。剩余差距的几种解释，即路线多样性、有偏奖励向锚点的牵引以及训练长度的差异，需要通过受控运行加以区分。表还可以支持测试时的搜索，或为元启发式提供热启动。上界只要求一个解的架次对其他顺序仍然可行，但它在其他问题上是否有用还取决于它有多宽松，这一点我们无法事先预测。多个载体 \citep{fahradyan2020} 和一艘船上的多架无人机 \citep{ren2026} 会带来不同的问题，因为此时架次可能在时间上重叠，航段结构也会改变。}

\section*{CRediT authorship contribution statement\quad CRediT 作者贡献声明}
\placeholder{To be completed by the authors.}

\section*{Declaration of competing interest\quad 利益冲突声明}
\placeholder{To be completed by the authors.}

\section*{Funding\quad 资助}
\placeholder{To be completed by the authors.}

\section*{Data availability\quad 数据可用性}
\placeholder{State where the code, generated instances and per-instance results are available. Real529 instances are from the cited sources.}

\section*{Ethics statement\quad 伦理声明}
This study used computational experiments on synthetic and published benchmark instances and involved no human participants or personal data.

\section*{Declaration of generative AI use\quad 生成式人工智能使用声明}
\placeholder{The authors should confirm this statement.} During the preparation of this work the authors used Claude (Anthropic) to assist with literature search and manuscript drafting. The authors reviewed and edited the content and take full responsibility for the publication.

\bibliographystyle{elsarticle-harv}
\bibliography{refs}

\end{document}
